\documentclass[10pt,a4paper]{amsart}
\usepackage{amsmath,amssymb,amsthm,graphicx}
\usepackage{hyperref}
\usepackage{amsmath}
\usepackage{amssymb}
\usepackage{amsthm}
\usepackage{amscd}
\usepackage{xcolor}

\usepackage{nicefrac, faktor}
\usepackage[shortlabels]{enumitem}
\usepackage{mathrsfs}

\numberwithin{equation}{section}

\newtheorem{Theorem}{Theorem}[section]
\newtheorem*{Theorem*}{Theorem}
\newtheorem{Corollary}[Theorem]{Corollary}
\newtheorem{Lemma}[Theorem]{Lemma}
\newtheorem{Proposition}[Theorem]{Proposition}
\newtheorem{Assertion}[Theorem]{Assertion}
\newtheorem{Conjecture}[Theorem]{Conjecture}
 { \theoremstyle{definition}
\newtheorem{Definition}[Theorem]{Definition}
\newtheorem{Note}[Theorem]{Note}
\newtheorem{Example}[Theorem]{Example}
\newtheorem{Remark}[Theorem]{Remark}
}

\newcommand{\F}{\mathbb{F}}
\newcommand{\Z}{\mathbb{Z}}
\newcommand{\R}{\mathbb{R}}
\newcommand{\C}{\mathbb{C}}
\newcommand{\N}{\mathbb{N}}
\newcommand{\SSS}{\mathbb{S}}
\newcommand{\TTT}{\mathbb{T}}
\newcommand{\rk}{\mathrm{rk}}
\newcommand{\Hom}{\mathrm{Hom}}
\newcommand{\SO}{\mathrm{SO}}
\newcommand{\charac}{\mathrm{char}}
\newcommand{\eps}{\varepsilon}
\newcommand{\Sc}{{\rm Sc}}
\newcommand{\scal}{\rm scal}
\newcommand{\SU}{\rm SU}
\newcommand{\id}{{\rm id}}
\newcommand{\MSO}{{\rm MSO}}
\newcommand{\sB}{{\sf B}}
\newcommand{\Tor}{{\rm Tor}}
\newcommand{\im}{{\rm im}}
\newcommand{\mmod}{{\rm mod}}
\newcommand{\Sq}{{\rm Sq}}
\newcommand{\lf}{{\rm lf}}

\newcommand{\myicon}{$\,\,\,\triangleright$}

\newcommand{\no}{{\color{brown!70!black}$\nexists$}{\sf PSC}}
\newcommand{\yes}{{\color {blue!70!black}$\exists$}{\sf PSC}}
\newcommand{\PSC}{{\sf PSC}}

\newcommand{\Ho}{\mathrm{H}}
\newcommand{\Be}{\sf{B}}


\AddToHook{cmd/maketitle/after}{\markboth{M. Gromov and B. Hanke}{Non-spin scalar curvature invariance}}
\title{Positive scalar curvature and homological invariance for non-spin universal covers}
\author{Misha Gromov, Bernhard Hanke}
\date{doi:10.3842/SIGMA.2024.069; CC BY 4.0}
\begin{document}
\maketitle
\noindent\textit{Allowed standard prerequisites (editorial).}\begin{itemize}
\item Author-cited wild nonabelian Hodge correspondence and harmonic bundle prerequisites. Original bibliography retained. All source-local filtered extension constructions and classification proofs retained.
\end{itemize}\medskip
\section{Introduction}

In his construction of counterexamples to the Gromov--Lawson--Rosenberg conjecture, Thomas Schick in \cite{Sch98} discovered a purely homological obstruction to the existence of metrics with positive scalar curvature on closed oriented smooth manifolds.
We say that a smooth manifold is~\yes, if it admits a Riemannian metric of positive scalar curvature, and \no, if it does not admit such a metric.
Schick's construction is based on the observation that by the descent method of Schoen--Yau \cite{S-Y1979,SY}, the \yes-property of a closed oriented manifold $V$ of dimension~${n \leq 7}$ implies that the Poincar\'e dual of a class $c \in H^{n-\ell}(V; \Z)$ which is the product of $1$-dimensional integral cohomology classes can be represented by a~closed oriented smoothly embedded \yes-submanifold $\Sigma \subset V$ of dimension $\ell$.
Now, if $\ell = 2$ and the fundamental class of $\Sigma$ can be detected by a product of two $1$-dimensional cohomology classes in $V$ with coefficients in some field $\F$, then $\Sigma$ is \no~by the Gauss--Bonnet theorem.
This cohomological condition on the fundamental class of $\Sigma$ is satisfied if there exist classes $c_1, c_2 \in H^1(V; \F)$ such that $c \cup c_1 \cup c_2 \in H^n(V; \F)$ evaluates to nonzero on the fundamental class of $V$.
It follows by contraposition that in this case $V$ is \no.

In the present paper, we generalize these ``toral'' \yes-obstructions in two directions.
First, we combine the torality condition on the Poincar\'e dual of $c$ with an enlargeability condition on $V$, see Definition \ref{def:enlargeable}.
The resulting largeness property depends only on the image of the fundamental class of $V$ in $H_n( \pi_1(V) ; \F)$ under the classifying map, see Proposition \ref{homol_inv}.
Second, we allow the case $\dim \Sigma = 3$ by exploiting the restricted nature of \yes-$3$-manifolds, see Proposition \ref{prop:toralthree}.
In Theorem \ref{thm:tornexist}, we obtain a general \yes-obstruction in this setting.

As an application, in Theorem \ref{thm:plenty}, we construct new examples of \no-manifolds whose Cartesian products are \yes, thus challenging our geometric intuition that \no-manifolds are ``large'' while \yes-manifolds are ``small''.


Some fundamental questions remain unanswered, most notably whether there exist toral \yes-manifolds of dimension $n \geq 4$ and with fundamental group $(\Z/p)^n$ for some odd prime $p$, see Conjecture \ref{main_conj} and \cite[Problem 9]{MP2021}.
We hope that this note will stimulate some interest in this and related problems, some of which we address in Section \ref{fourconjectures}.

In an appendix, we prove homological invariance of the \yes-property of closed oriented manifolds of dimension at least $5$ with non-spin universal covers, filling a gap in an earlier treatment in the literature \cite{RS}.
This result is of independent interest.

\section[Torality and tor-enlargeability]{$\boldsymbol{\F}$-torality and tor-$\boldsymbol{\ell/\F}$-enlargeability}

\begin{Definition} \label{deftoral}
Let $X$ be a topological space, let $\ell$ be a non-negative integer and let $\mathbb{F}$ be a~field, e.g., $\mathbb{Q}$ or the finite field $\mathbb F_{p}$ for some prime $p$.

In the line of \cite{Han20,Jo2004}, a homology class $h \in H_\ell(X; \F)$ is called {\it $\mathbb F$-toral} if it is detected by a~product of $1$-dimensional cohomology classes, i.e., there exist classes $c_1, \ldots, c_\ell \in H^1(X ; \mathbb F)$ such that
\[
 \langle c_1\cup \cdots \cup c_\ell, h \rangle \neq 0 \in \F.
\]
Here we use the Kronecker pairing (evaluation pairing) of cohomology and homology.
We agree that all $h \neq 0 \in H_0(X; \F)$ are $\F$-toral.
Otherwise, the class $h$ is called {\it $\F$-atoral}.
Obviously, the subset of $\F$-atoral classes forms a linear subspace of $H_{\ell}(X; \F)$.

A closed, i.e., compact-without-boundary, $\F$-orientable\footnote{This condition is empty for $\charac(\F) = 2$ and equivalent to orientability for $\charac(\F) \neq 2$.} $n$-dimensional manifold $V$ is called {\it $\F$-toral}, if, after choosing some $\F$-orientation, its fundamental class $[V]_{\F} \in H_{n}(V; \F)$ is $\F$-toral.
Otherwise, $V$ is called {\it $\F$-atoral.}
\end{Definition}

In the following, let $\SSS^n$ denote the unit $n$-sphere equipped with the standard round metric, and let $\infty \in \SSS^n$ be some base point.
We call a continuous map $f \colon X \to \SSS^n$ {\em compactly supported}, if there exists a compact subset $K \subset X$ with $f(X \setminus K) = \{ \infty\}$.

\begin{Remark} \label{rem:bemerkung}\quad
\begin{enumerate}[(a)]\itemsep=0pt
\item If $\charac(\F) \neq 2$, then the classes $c_1, \ldots, c_\ell$ appearing in Definition \ref{deftoral} span an exterior algebra on $\ell$ generators of dimension $1$ in the $\F$-algebra $H^*(X ; \F)$.
\item \label{overtorical} A closed oriented $n$-manifold $V$ is $\mathbb{Q}$-toral, if and only if it admits a map of non-zero degree to the $n$-torus $\TTT^n$.
In other words, $V$ is {\em overtorical}.
This follows from the universal coefficient formula $H^*(V; \mathbb{Q}) = H^*(V; \mathbb{Z}) \otimes \mathbb{Q}$ and the one-to-one correspondence of classes in $H^1(V; \mathbb{Z})$ and homotopy classes of maps $V \to S^1$.
\end{enumerate}
\end{Remark}


\begin{Definition} \label{def:enlargeable} Let $n \geq 1$ and $0 \leq \ell \leq n$.
An $\F$-orientable Riemannian $n$-manifold $(V,g)$ is called {\it tor-$\ell/\mathbb F$-enlargeable}, if for all $\varepsilon > 0$, there exists a Riemannian covering
$
 \bar V = \bar V_\varepsilon \to V
$
and a compactly supported smooth $\varepsilon$-Lipschitz map
$
 f = f_\varepsilon \colon \bar V \to \SSS^{n-\ell}
$,
which is of {\em non-zero $\F$-toral degree}.
By definition, this means that for some, hence every, regular value ${s \in \SSS^{n - \ell} \setminus \{\infty\}}$ of~$f$, the homology class $\big[f^{-1}(s)\big] \in H_{\ell}\big(\bar V; \F\big)$ represented by the closed $\F$-orientable submanifold~${f^{-1}(s) \subset \bar V}$ is $\F$-toral.

For $\ell = n$, this means that a finite cover of $V$ is $\F$-toral.
\end{Definition}

\begin{Remark}\quad
\begin{enumerate}[(a)]\itemsep=0pt
\item In the case of closed $V$, this property is independent of the choice of $g$.
Thus, it is an invariant of the diffeomorphism type of $V$.
\item In fact, if $V$ is closed, connected and $\F$-oriented, this property depends only on the image of the fundamental class $[V]_\F \in H_n(V; \F)$ under the classifying map $V \to {\sf B}\pi_1(V)$, see Proposition \ref{homol_inv}.
In particular, for closed $V$, tor-$\ell/\F$-enlargeability is an invariant of the homotopy type of $V$.
\item Tor-$0/ \mathbb{Q}$-enlargeable $n$-manifolds $V$ are {\em enlargeable} in the classical sense of \cite[Section 5]{G-L1983}, i.e., for all $\eps > 0$, there exists a Riemannian covering $\bar V \to V$ together with a compactly supported $\eps$-Lipschitz map $\bar V \to \SSS^n$ of non-zero degree.
\item ``Tor-$\ell/\F$-enlargeable'' generalizes ``$\mathcal{SYS}$-enlargeable'' from \cite{Gr2018}, where $\mathcal{SYS}$ stands for \textbf Schoen--\textbf Yau--\textbf Schick.
\end{enumerate}
\end{Remark}

\begin{Example} \label{ex:toral}\quad
\begin{enumerate}[(a)]\itemsep=0pt
\item \label{eins} For all $n \geq 1$ and all $\F$, the $n$-torus $V = \TTT^n = S^1 \times \cdots \times S^1$ is $\F$-toral.
\item \label{projspace} For all $n \geq 1$, the real projective space $\mathbb{R} P^n$ is $\F_2$-toral.
\item \label{ex:overtorical} Let $0 \leq \ell \leq n$ and let $V$ be a closed $\F$-orientable $n$-manifold together with a map $f \colon V \to {\TTT}^{n-\ell}$ of non-zero $\F$-toral degree.
Then $V$ is $\F$-toral and tor-$\ell/\F$-enlargeable.
Such $V$ are called {\em tor-$\ell/\F$-overtorical}.
\item Let $V$ be tor-$\ell / \F$-enlargeable and let $W$ be a an $\F$-orientable Riemannian manifold admitting a proper $1$-Lipschitz map $W \to V$ of non-zero $\F$-toral degree.
Then $W$ is tor-$\ell/\F$-enlargeable.
\end{enumerate}
\end{Example}


For a group $G$, let ${\sf B}G$ be the classifying space of $G$ and let $H_*(G) = H_*({\sf B}G)$ be the group homology of $G$.
For a connected topological space $X$, let
$\varphi \colon  X \to {\sf B} \pi_1(X)$
be the classifying map of the universal cover of~$X$.
It is unique up to homotopy.

\begin{Definition} \label{def:essential} Let $V$ be a closed connected manifold of dimension $n$.
\begin{enumerate}[(a)]\itemsep=0pt
\item Let $\F$ be a field.
$V$ is called {\em $\F$-essential} if it is $\F$-orientable and the classifying map $\varphi \colon V \to {\sf B}\pi_1(V)$ sends some fundamental class $[V] _\F \in H_{n}(V; \F)$ to a non-zero homology class in $H_n(\pi_1(V); \F)$.
\item $V$ is called {\em $($integrally$)$ essential} if it is orientable and $\varphi$ sends some integral fundamental class $[V] \in H_n(V;\Z)$ to a non-zero homology class in $H_n(\pi_1(V); \Z)$.
\end{enumerate}
\end{Definition}

Any oriented $\F$-essential manifold $V^n$ is essential, since $\varphi_*([V]_\F) \in H_n( \pi_1(V) ; \F)$ is just the $\F$-reduction of $\varphi_*([V]) \in H_n(\pi_1(V) ; \Z)$.
However, the last class may be a {\em torsion class}, see Example \ref{ex:torsion}.
Hence, the title of our paper.

In the following, for $n \geq 0$, let $\sigma^n \in H^{n}_{\rm c}(\SSS^{n} \setminus \{ \infty\}; \F)$ be the non-zero compactly supported cohomology class Poincar\'e dual to the class in $H_0( \SSS^{n} \setminus \{ \infty\}; \F) = \F$ represented by a point.
Note that for each compactly supported map $f \colon X \to \SSS^{n}$, the cohomology class $f^*(\sigma^n) \in H^{n}_{\rm c}(X; \F)$ is compactly supported.

\begin{Proposition} \label{special} Each connected tor-$\ell / \F$-enlargeable manifold $V$ is $\F$-essential.
\end{Proposition}

\begin{proof}
We construct a CW model of ${\sf B}\pi_1(V)$ by attaching cells of dimension $\geq 3$ to $V$, killing all homotopy groups in dimension $\geq 2$.
In this description, the classifying map $\varphi \colon V \to {\sf B}\pi_1(V)$ becomes an inclusion.

Let $n = \dim V$ and assume that $\varphi_*([V]_\F) = 0 \in H_n ( \pi_1(V) ; \F)$.
Then there exists a connected finite subcomplex $V \subset S \subset {\sf B}\pi_1(V)$ such that the inclusion map $\iota \colon V \hookrightarrow S$ satisfies $\iota_*([V]_{\F}) =0 \in H_n(S; \F)$.
Since $V$ and $S$ have the same $2$-skeleton, the map $\iota$ induces an isomorphism of fundamental groups.

By assumption and Poincar\'e duality, for all $\varepsilon > 0$, there exists a connected covering $\pi_V \colon \bar V = \bar V_\eps \to V$, a compactly supported smooth $\eps$-Lipschitz map $f_\eps \colon \bar V \to \SSS^{n-\ell}$, and classes $c_1, \ldots , c_{\ell} \in H^1\big(\bar V; \F\big)$ so that
\begin{equation} \label{thetadeg}
 \big( f_\eps^*\big(\sigma^{n-\ell}\big) \cup (c_1 \cup \cdots \cup c_\ell ) \big) \cap \pi_V^!( [V]_{\F}) \neq 0 \in H_0 \big( \bar V; \F\big) .
\end{equation}
 Here $H_*^{\rm lf}$ is locally finite homology and $\pi_V^{!} \colon H_n(V ; \F) \to H^{\rm lf}_n\big( \bar V; \F\big)$ is the homological transfer.
Note that $f_\eps^*\big(\sigma^{n-\ell}\big) \cup (c_1 \cup \cdots \cup c_\ell ) \in H_c^n\big( \bar V; \F\big)$, so that the class considered in~\eqref{thetadeg} indeed lies in $H_0\big(\bar V; \F\big)$.

 Let $G := \pi_1\big(\bar V\big) < \pi_1(V)$ and let $\pi \colon {\sf B}G \to {\sf B}\pi_1(V)$ be the corresponding covering map.
Let~${\bar S := \pi^{-1}(S) \subset {\sf B}G}$.
Since $S \hookrightarrow {\sf B}\pi_1(V)$ induces an isomorphism of fundamental groups, $\bar S$ is connected and the inclusion $\bar V \hookrightarrow \bar S$ induces an isomorphism of fundamental groups.
Hence, we obtain $H^1\big(\bar S; \F\big) \cong H^1\big(\bar V; \F\big)$, and the classes $c_1, \ldots, c_\ell \in H^1\big(\bar V; \F\big)$ extend to classes $c_1', \ldots, c_\ell' \in H^1\big(\bar S; \F\big)$.

By \cite[Lemma 3.2]{BH10} and since $S$ is a finite complex, one can choose $\eps > 0$ at the beginning of this argument so small that the map $f_\eps \colon \bar V \to \SSS^{n-\ell}$ extends to a compactly supported map $F \colon \bar S \to \SSS^{n-\ell}$.

Let $\pi_S^! \colon H_n(S; \F) \to H^{\rm lf}_n\big(\bar S; \F\big)$ be the homological transfer map for the covering $\pi_S \colon \bar S \to S$.
Then \eqref{thetadeg} together with the naturality of the transfer map implies
\[
 \big( F^*\big(\sigma^{n-\ell}\big) \cup ( c_1' \cup \cdots \cup c_\ell') \big) \cap \pi_S^!(\iota_*( [V]_{\F})) \neq 0 \in H_0\big(\bar S; \F\big) .
\]
This contradicts our assumption $\iota_*([V]_\F) = 0 \in H_n(S; \F)$.
\end{proof}

Extending this line of thought, in Proposition \ref{homol_inv}, we will show a homological invariance property of tor-$\ell / \F$-enlargability.
This is similar in spirit to \cite[Section 3]{BH10}.
In the following, we equip connected simplicial complexes with their canonical path metrics whose restriction to each $k$-simplex is the standard metric on $\Delta^k \subset \R^{k+1}$.

\begin{Definition} \label{largehom}
Let $0 \leq \ell \leq n$.
Let $X$ be a connected simplicial complex with finitely presented fundamental group.
A homology class $h \in H_n(X; \F)$ is called {\em tor-$\ell / \F$-enleargeable}, if there exists a finite connected subcomplex $S \subset X$ and a class $h_S \in H_n(S; \F)$ such that
\begin{enumerate}[(a)]\itemsep=0pt
\item \label{isofund} The inclusion $S \hookrightarrow X$ induces an isomorphism of fundamental groups\footnote{We implicitly assume that $S$ contains the base point of $X$.} and the induced map $H_n(S ; \F) \to H_n(X; \F)$ sends $h_S$ to $h$;
\item \label{shrinkingmap} for all $\eps > 0$, there exists a connected cover $\pi_S \colon \bar S = \bar S_\eps \to S$, a compactly supported $\eps$-Lipschitz map $f_\eps \colon \bar S \to \SSS^{n-\ell}$ and classes $c_1 , \ldots, c_\ell \in H^1\big(\bar S; \F\big)$ that satisfy
 \begin{equation*}
 \big( f_\eps^*\big(\sigma^{n-\ell}\big) \cup (c_1 \cup \cdots \cup c_\ell) \big) \cap \pi_S^!(h_S) \neq 0 \in H_{0}\big(\bar S; \F\big) .
\end{equation*}
\end{enumerate}
Let
\smash{$
 H_*^{\text{\rm tor-}\ell/\F\text{\rm-enl}}(X; \F) \subset H_*(X ; \F)
$}
be the subset of tor-$\ell/\F$-enlargeable homology classes.
\end{Definition}

The complex $S$ in Definition \ref{largehom} can be made arbitrarily large:

\begin{Lemma} \label{subcomplexinv} Let $h\in H_n(X; \F)$ be tor-$\ell / \F$-enlargeable.
Let $S \subset X$ and $h_S \in H_n(S;\F)$ as in Definition {\rm\ref{largehom}} and let $K \subset X$ be a finite subcomplex.

Then there exists a finite connected subcomplex $S \cup K \subset T \subset X$ such that the inclusion~${T \hookrightarrow X}$ induces an isomorphism of fundamental groups.

Given such a subcomplex $T \subset X$, let $h_T \in H_n(T; \F)$ be the image of $h_S$ under the inclusion~${S \hookrightarrow T}$.
Then $T$ and $h_T$ satisfy properties {\rm\ref{isofund}} and {\rm\ref{shrinkingmap}} in Definition {\rm\ref{largehom}}.
\end{Lemma}

\begin{proof} The existence of $T$ is clear.
Obviously, the induced map $H_n(T; \F) \to H_n(X; \F)$ sends~$h_T$ to $h$.
Let $\eps > 0$, let $f_{\eps} \colon \bar S \to \SSS^{n-\ell}$ be an $\eps$-contracting compactly supported map where $\bar S \to S$ is a connected cover and let $c_1, \ldots, c_\ell \in H^1\big(\bar S; \F\big)$ such that
\begin{equation} \label{nonzero}
 \big( f_\eps^*\big(\sigma^{n-\ell}\big) \cup (c_1 \cup \cdots \cup c_\ell) \big) \cap \pi_S^!(h_S) \neq 0 \in H_{0}\big(\bar S; \F\big) .
\end{equation}
Let $\bar T \to T$ be the uniquely determined connected cover which restricts to $\bar S \to S$.
The inclusion~${\iota\colon \bar S \hookrightarrow \bar T}$ induces an isomorphism of fundamental groups, hence an isomorphism $H^1\big(\bar T; \F\big) \cong H^1\big(\bar S; \F\big)$.

We claim that, for small enough $\eps$, the map $f_\eps$ extends to a compactly supported $(C \cdot \eps)$-contracting map $F_{C\eps} \colon \bar T \to \SSS^{n-\ell}$ where $C > 0$ depends only on $S$ and $T$, but not on $\eps$.

The proof is by induction on the $k$-skeleta $T^{(k)} \subset T$ relative to $S$, where $k \in \N_0$.
If $S^{(0)} = T^{(0)}$, this is achieved by applying \cite[Lemma 3.2]{BH10} finitely many times.
If $S^{(0)} \subsetneq T^{(0)}$, the proof is a~bit more complicated as the cases $k=0$ and $k=1$ have to be treated simultaneously.
We refer the reader to the detailed argument in \cite[p.\ 474f.]{BH10}, which carries over to our case.

Let $c_1', \ldots c_\ell' \in H^1\big(\bar T; \F\big) \cong H^1\big(\bar S; \F\big)$ correspond to $c_1, \ldots, c_\ell \in H^1\big(\bar S; \F\big)$.
Using the naturality of the transfer map, \eqref{nonzero} implies
\begin{equation*}
 \big( (F_{C \eps})^*\big(\sigma^{n-\ell}\big) \cup (c'_1 \cup \cdots \cup c'_\ell) \big) \cap \pi_T^!(h_T) \neq 0 \in H_0\big(\bar T; \F\big).
\end{equation*}

Since $\eps > 0$ can be chosen arbitrarily small, this implies property \ref{shrinkingmap} of Definition \ref{largehom} for the class $h_T \in H_n(T; \F)$.
\end{proof}

\begin{Corollary} The complement $H_*(X ; \F) \setminus H_*^{\text{\rm tor-}\ell/\F\text{\rm -enl}}(X; \F) \subset H_*(X ; \F)$ is a linear subspace.
\end{Corollary}

\begin{proof} The most important part is to show that if $h_1 , h_2 \in H_*(X; \F)$ such that $h_1 + h_2 \in \smash{H_*^{\text{\rm tor-}\ell/\F\text{\rm -enl}}(X; \F)}$, then at least one of $h_1$ or $h_2$ is contained in \smash{$H_*^{\text{\rm tor-}\ell/\F\text{\rm -enl}}(X; \F)$}.
This follows because for the subcomplex $S \subset X$ witnessing \smash{$h_1 + h_2 \in H_*^{\text{\rm tor-}\ell/\F\text{\rm -enl}}(X; \F)$} we can assume by Lemma \ref{subcomplexinv} that both $h_1$ and $h_2$ lie in the image of $H_*(S; \F) \to H_*(X; \F)$.
The remaining details are left to the reader.
\end{proof}

\begin{Lemma} \label{invariance} Let $X$ and $X'$ be connected simplicial complexes with finitely presented fundamental groups.
Let $\eta \colon X \to X'$ be a map inducing an isomorphism of fundamental groups.
Then
\[
 (\eta_*)^{-1} \big( H_*^{\text{\rm tor-}\ell/\F\text{\rm -enl}}(X'; \F)\big) = H_*^{\text{\rm tor-}\ell/\F\text{\rm -enl}}(X; \F) .
\]
\end{Lemma}

\begin{proof}
We first show the inclusion ``$\subset$''.
This also implies the inverse inclusion, if $\eta$ is a~homotopy equivalence, by considering a~homotopy inverse of $\eta$.
Suppose $h \in H_*(X; \F)$ such that~\smash{$h' := \eta_*(h) \in H_*^{\text{\rm tor-}\ell/\F\text{\rm -enl}}(X'; \F)$}.
We need to show that \smash{$h \in H_*^{\text{\rm tor-}\ell/\F\text{\rm -enl}}(X; \F)$}.


There exists a finite connected subcomplex $S \subset X$ such that the inclusion $S \hookrightarrow X$ induces an isomorphism of fundamental groups and a class $h_S \in H_n(S; \F)$ which maps to $h \in H_n(X; \F)$.

Using Lemma \ref{subcomplexinv}, there exists a finite connected subcomplex $S' \subset X$ and $h'_{S'} \in H_n(S'; \F)$ satisfying the properties \ref{isofund} and \ref{shrinkingmap} of Definition \ref{largehom} for $h'$ such that, in addition, $\eta(S) \subset S'$ and~${\eta_*(h_S) = h'_{S'}}$.
Since $S$ is compact, the restriction $\eta|_{S} \colon S \to S'$ is, up to homotopy, $\lambda$-Lipschitz for some $\lambda > 0$.

Let $\eps > 0$, let $c'_1, \ldots, c'_\ell \in H^1\big(\bar S' ; \F\big)$, let $\pi_{S'} \colon \bar S' \to S'$ be a connected cover and let $f \colon \bar S' \to \SSS^{n-\ell}$ be a compactly supported $\tfrac{\eps}{\lambda}$-contracting map satisfying
\begin{equation} \label{forSprime}
 \big( f^*\big(\sigma^{n-\ell}\big) \cup (c'_1 \cup \cdots \cup c'_\ell) \big) \cap \pi_{S'}^!(h'_{S'}) \neq 0 \in H_{0}\big(\bar S' ; \F\big) .
\end{equation}
Since $\eta$ induces an isomorphism of fundamental groups, there is a unique connected cover $\pi_S \colon \bar S \to S$ such that $\eta$ lifts to a proper map $\bar \eta \colon \bar S \to \bar S'$ which induces an isomorphism of fundamental groups, hence an isomorphism $H^1\big(\bar S' ; \F\big) \cong H^1\big(\bar S; \F\big)$.
The lift $\bar \eta$ is still $\lambda$-Lipschitz continuous (recall that $\bar S$ is equipped with the canonical path metric).
Then $f \circ \bar \eta \colon \bar S \to \SSS^{n-\ell}$ is $\eps$-Lipschitz and compactly supported since $\bar \eta$ is proper.


Let $c_1 , \ldots, c_\ell \in H^1\big(\bar S; \F\big)$ correspond to $c'_1, \ldots, c'_\ell \!\in\! H^1\big(\bar S' ; \F\big)$.
By \eqref{forSprime} and since ${\eta_*(h_S)\! =\! h'_{S'}}$, we get
\begin{equation*}
 \bar \eta_* \big( \big( (f \circ \bar \eta)^*\big(\sigma^{n-\ell}\big) \cup (c_1 \cup \cdots \cup c_\ell) \big) \cap \pi_S^!(h_S) \big) \neq 0 \in H_0\big(\bar S' ; \F\big) .
 \end{equation*}
Since $\bar \eta_* \colon H_0\big( \bar S; \F\big) \to H_0\big(\bar S' ; \F\big)$ is an isomorphism, this shows that property \ref{shrinkingmap} of Definition~\ref{largehom} holds for $(S, h_S)$.
We conclude \smash{$h \in H_*^{\text{\rm tor-}\ell/\F\text{\rm -enl}}(X; \F)$}, which was our claim.

To show the reverse inclusion ``$\supset$'' in general, after replacing $X'$ by a homotopy equivalent simplicial complex, we may assume that $\eta \colon X \to X'$ is a simplicial inclusion.
Let \smash{$h \in H_*^{\text{\rm tor-}\ell/\F\text{\rm -enl}}(X; \F)$} and let $S \subset X$ and $h_S \in H_n(S; \F)$ be such that the properties \ref{isofund} and \ref{shrinkingmap} of Definition \ref{largehom} hold for $(S, h_S)$.
Then $S \hookrightarrow X'$ induces an isomorphism of fundamental groups and $H_n(S; \F) \to H_n(X'; \F)$ sends $h_S$ to $\eta_*(h)$.
This shows \smash{$\eta_*(h) \in H_*^{\text{\rm tor-}\ell/\F\text{\rm -enl}}(X'; \F)$}, finishing the proof of the reverse inclusion.
\end{proof}

For each finitely presented group $G$, we now define a subset
\[
 H_*^{\text{\rm tor-}\ell/\F\text{\rm -enl}}(G; \F) \subset H_*(G; \F)
\]
of {\em tor-$\ell / \F$-enleargeable classes} in the group homology of $G$ by setting
\[
 H_*^{\text{\rm tor-}\ell/\F\text{\rm -enl}}(G; \F) := H_*^{\text{\rm tor-}\ell/\F\text{\rm -enl}}(X; \F),
\]
where $X$ is an arbitrary connected simplicial complex representing $\sB G$.
By Lemma \ref{invariance}, this definition is independent of the choice of $X$.

\begin{Proposition} \label{homol_inv} A connected closed $\F$-oriented manifold $V$ is tor-$\ell/\F$ enlargeable if and only if the image of the fundamental class $[V]_{\F} \in H_n(V; \F)$ under the classifying map $V \to {\sf B}\pi_1(V)$ lies in \smash{$H_*^{\text{\rm tor-}\ell / \F\text{\rm -enl}}(\pi_1(V); \F)$}.
\end{Proposition}

\begin{proof} Choose a triangulation of $V$.
Then $V$ is tor-$\ell/\F$-enlargeable if and only if the fundamental class $[V]_\F \in H_n(V ; \F)$ is tor-$\ell/\F$-enlargeable.
This holds because the path metric induced by any Riemannian metric on $V$ is bi-Lipschitz equivalent to the simplicial path metric for the chosen triangulation on $V$ and each compactly supported $\eps$-Lipschitz map $\bar V \to \SSS^{n}$ of non-zero $\F$-toral degree is homotopic to a smooth compactly supported $2\eps$-Lipschitz map~${\bar V \to \SSS^{n-\ell}}$ of non-zero $\F$-toral degree.
Furthermore, the classifying map $V \to \sB \pi_1(V)$ induces an isomorphism of fundamental groups.
Thus the claim follows from Lemma \ref{invariance} for $X = V$ and~${X' = \sB \pi_1(V)}$.
\end{proof}



\begin{Example} \label{ex:torsion} Let $p \geq 2$ be a prime, let $\ell \geq 1$ and let $\tau_\ell \in H_\ell \big( (\Z/p)^\ell ; \F_p\big)$ be represented by the map
$
 \TTT^\ell = \sB \Z^{\ell} \to \sB (\Z/p)^{\ell}$
which is induced by the canonical projection $\Z^{\ell} \to (\Z/p)^{\ell}$.
Apparently, $\tau_\ell \in H_\ell\big( (\Z/p)^{\ell} ; \F_p\big)$ is $\F_p$-toral.

Next, let $(W_0, g)$ be a closed connected oriented Riemannian manifold of dimension $n$ where~$g$ has non-positive sectional curvature.
Let $\tilde W_0 \to W_0$ be the universal cover.
By the Cartan--Hadamard theorem, for any $x \in \tilde W_0$, the exponential map $T_x \tilde W_0 \to \tilde W_0$ is a diffeomorphism with $1$-Lipschitz inverse with respect to the lifted metric on $\tilde W_0$.
In particular, $W_0$ is a model for~$\sB \pi_1(W_0)$, and for all $\eps>0$ there exists a compactly supported $\eps$-contracting map $\tilde W_0 \to \SSS^n$ of degree $1$.
We claim that \smash{$\tau_\ell \times [W_0] \in H_{\ell + n}^{\text{\rm tor-}\ell/{\F_p}\text{\rm -enl}}\big((\Z/p)^{\ell} \times \pi_1(W_0) ; \F_p\big)$}.

To prove this, choose a simplicial complex $X$ representing $\sB (\Z/p)^{\ell}$.
Since $(\Z/p)^{\ell}$ is a finite group, one can choose $X$ finite in each dimension.
Let $K = X^{(2\ell)} \subset X$ be the $2\ell$-skeleton of $X$.
The inclusion $K \hookrightarrow X$ is $2 \ell $-connected.
In particular there is a unique $h_K \in H_\ell( K ; \F_p)$ mapping to $\tau_\ell$ under the map $H_*(K ; \F_p) \to H_*(X ; \F_p)$.


Choose some triangulation of $W_0$ and notice that the canonical path metric on $W_0$ is bi-Lipschitz equivalent to the metric induced by $g$.
Equip $X \times W_0$ with the product triangulation.
Then $K \times W_0$ with its product triangulation is a finite connected subcomplex of $X \times W_0$ and the class $h_K \times [W_0] \in H_{\ell + n}(K \times W_0 ; \F_p)$ maps to $\tau_\ell \times [W_0]$.
Furthermore, the inclusion $K \times W_0 \hookrightarrow X \times W_0$ induces an isomorphism of fundamental groups.


Let $\eps> 0 $, and let $f \colon \tilde W_0 \to \SSS^{n}$ be a compactly supported $\eps$-contracting map of degree $1$.
Then the composition
\[
 K \times \tilde W_0 \to \tilde W_0 \to \SSS^{n}
\]
is compactly supported, $\eps$-contracting (the first map is $1$-Lipschitz) and of non-zero $\F_p$-toral degree, as required.

As an application, assume that $\ell + n \geq 4$ and let an oriented manifold $V$ be obtained by surgery along $\ell$ pairwise disjoint embedded circles in $\TTT^{\ell} \times W_0$ representing the $p$-multiples of generators of the fundamental group $\pi_1\big({\TTT}^\ell\big) = \Z^\ell \subset \pi_1\big(\TTT^{\ell} \times W_0\big)$.


Then $\pi_1(V) \cong (\Z/p)^\ell \times \pi_1(W_0)$ and the classifying map $\varphi \colon V \to \sB \pi_1(V)$ is oriented bordant over $\sB \pi_1(V)$ to the map $\TTT^{\ell} \times W_0 = \sB \Z^{\ell} \times W_0 \to \sB (\Z/p)^{\ell} \times \sB \pi_1(W_0) = \sB \pi_1(V)$.
In particular, we obtain $\varphi_{*}([V]_{\F_p}) = \tau_\ell \times [W_0]$.
By Proposition \ref{homol_inv}, the manifold $V$ is tor-$\ell/\F_p$-enlargeable.

Note that $\varphi_*([V]) \in H_{\ell+n}(\pi_1(V); \Z)$ is a $p$-torsion class.
\end{Example}

\begin{Remark} One may define a quantitative version of tor-$\ell/\F$-enlargeability based on the notion of {\em $\tilde \square^{\perp}$-spread} of $\F$-toral homology classes rather than enlargeability, compare \cite[Section~7.1]{Gr21}.
\end{Remark}


\section[Atorality of PSC-3-manifolds]{Atorality of \yes-$\boldsymbol{3}$-manifolds} \label{nopsc}

We will show the following homological property of \yes-$3$-manifolds.

 \begin{Proposition} \label{prop:toralthree}
Let $\F$ be a field with $\charac( \F) \neq 2$.
Then closed orientable \yes-$3$-manifolds are $\F$-atoral.
\end{Proposition}

For the proof, we recall the following classification result of \yes-$3$-manifolds \cite{G-L1983}.

\begin{Proposition} \label{prop:decompthree} Closed connected orientable \yes-$3$-manifolds are diffeomorphic to connected sums of manifolds of the following type:
\begin{enumerate}[\myicon]\itemsep=0pt
\item closed connected $3$-manifolds with universal covers homotopy equivalent to $\SSS^3$,
\item copies of $\SSS^2\times \SSS^1$.
\end{enumerate}
In particular, they are not $\mathbb{Q}$-essential.
\end{Proposition}

Indeed, by the Kneser--Milnor prime decomposition \cite{Milnor1962}, every closed connected orientable $3$-manifold $V$ is diffeomorphic to the connected sum of manifolds of the form described in Proposition \ref{prop:decompthree} and aspherical manifolds, i.e., manifolds with contractible universal covers.

It remains to show that if $V$ is \yes, then aspherical summands cannot occur.
This non-trivial fact can be proved in several ways, see
\begin{enumerate}[\myicon]\itemsep=0pt
 \item \cite[Theorem 5.2]{S-Y1979}, based on the theory of {\em minimal hypersurfaces}, together with the resolution of Waldhausen's surface subgroup conjecture \cite{K-M2012};
\item \cite[Theorem 8.1]{G-L1983}, using the {\em index theory of the Dirac operator} on a (possibly non-compact) cover of $V$;
\item \cite[Theorem 3.4]{Ros1986}, showing that the {\em Rosenberg index} $\alpha(V) \in {\rm KO}_3(C^*_{\R} (\pi_1(V)))$ vanishes, together with the validity of the real Baum--Connes conjecture for $3$-manifold groups \cite{MOP08}, which in turn follows from the hyperbolization theorem of Thurston--Perelman;
\item \cite[Corollary 3.10.1\,(F$''$)]{Gr21}, estimating the {\em Uryson width} of the universal cover of $V$ and using the fact that groups which are quasi-isometric to trees are virtually free, see \cite[Theorem~7.19]{GdlH}.\footnote{For a remarkable generalization of this argument to manifolds in dimensions~4 and~5, see \cite{Ch-Li}.}
\end{enumerate}

\begin{Proposition} \label{ranks}
Let $V$ be a closed connected $n$-manifold, $n \geq 2$, with universal cover homotopy equivalent to $\SSS^n$.
Then $\dim_{\F} H^1 (V; \F) \leq 2$ for any field $\F$.
\end{Proposition}

\begin{proof} Let $G = \pi_1(V)$ and write $V = \tilde V / G$ where $\tilde V \to V$ is the universal cover.
Note that~$\tilde V$ is closed and orientable.
If the (free) action of $G$ on $\tilde V \simeq \SSS^n$ is orientation reversing, then $n$ is even and~${G = \Z/2}$ by the Lefschetz fixed point theorem.
In this case, we have $\dim_{\F} H^1(V ; \F) = 1$ if~${\charac(\F) = 2}$ and $\dim_{\F} H^1(V ; \F) = 0$ if $\charac(\F) \neq 2$.

So we can assume that the $G$-action on $\tilde V$ is orientation preserving.
A transfer argument shows that $H^*(V; \mathbb{Q}) \cong H^*(\SSS^n ; \mathbb{Q})$ and it remains to study the case when $\F$ is of finite characteristic~${p \geq 2}$.
Let $P < G$ be a Sylow $p$-subgroup and let $\pi \colon \tilde V / P \to \tilde V / G$ be the induced covering map.

Since $\tilde V$ is a homology $n$-sphere, it follows from \cite[Section XVI.9, Application 4]{C-E1956} that $P$ has periodic group cohomology of period $n+1$.
Using the classification of finite $p$-groups with periodic cohomology in \cite[Theorem XII.11.6]{C-E1956},
we conclude that $P$ is a cyclic $p$-group or generalized quaternion group.\footnote{Also compare \cite[Chapter IV.6]{Adem-Milgram}.}
Hence, $\dim_{\F} H^1\big(\tilde V / P; \F\big) = \dim_{\F} H^1(P ; \F) \leq 2$.
In fact, if~${p \geq 3}$, then $P$ is cyclic (since generalized quaternion groups are of even order), and hence~\smash{$\dim_{\F} H^1\big(\tilde V / P ; \F\big) \leq 1$}.

The induced map $\pi^* \colon H^*\big(\tilde V / G; \F\big) \to H^*\big(\tilde V / P ; \F\big)$ is injective since its composition with the cohomological transfer map \smash{$H^*\big(\tilde V / P; \F\big) \to H^*\big(\tilde V / G; \F\big)$} is multiplication by $\deg \pi = [G:P]$, which is prime to $p$.

From this follows the assertion of the proposition.
 \end{proof}

 \begin{Corollary} \label{cor:asphthree} Let $\charac(\F) \neq 2$ and let $V$ be a closed connected $\F$-toral $3$-manifold.
Then $V$ is diffeomorphic to a connected sum $K \sharp V'$ where $K$ is an aspherical $3$-manifold.
In particular, $V$ is $\mathbb{Q}$-essential.
\end{Corollary}

 \begin{proof}
By Proposition \ref{ranks}, and since for $\charac (\F) \neq 2$, every class in $H^1(V; \F)$ squares to $0$, the prime decomposition of $V$ must contain aspherical summands.
\end{proof}

The proof of Proposition \ref{prop:toralthree} follows from the combination of Proposition \ref{prop:decompthree} and Corollary~\ref{cor:asphthree}.

\begin{Remark} \label{rem:mixed}\quad
\begin{enumerate}[(a)]\itemsep=0pt
\item \label{leqzwei} It is easy to see that for $n \leq 2$ and all fields $\F$, all orientable connected $ \F$-toral $n$-manifolds are aspherical, hence $\mathbb{Q}$-essential, and are \no.
The $\F_2$-toral \yes-manifold $\mathbb{R}P^2$ shows that the orientability assumption cannot be dropped if $\charac(\F) = 2$.
\item \label{notinge4} Let $p$ be a prime, let $n \geq 4$ and let the oriented manifold $V^n$ be obtained by surgery along~$n$ pairwise disjoint embedded circles in $\TTT^n$ which represent the $p$-multiples of generators of the fundamental group $\pi_1({\TTT}^n) = \Z^n$.
Then $\pi_1(V) \cong (\Z/p)^n$, $V$ is $\F_p$-toral, but $V$ is not $\mathbb{Q}$-essential.
This shows that Corollary \ref{cor:asphthree} does not hold for $\F_p$-toral $n$-manifolds if $n \geq 4$.
\item The $\F_2$-toral \yes-manifold $\mathbb{R} P^3$ shows that Proposition \ref{prop:toralthree} does not hold for $p= 2$.
\end{enumerate}
\end{Remark}

\section[Torsion PSC-obstructions in higher dimensions]{Torsion \yes-obstructions in higher dimensions}

Using generalized minimal hypersurfaces ($\mu$-bubbles) and a Schoen--Yau-Schick descent, we propagate the \yes-obstruction from Section \ref{nopsc} to dimensions $\geq 4$.

We start with the following version of an overtorical band width inequality.
In the following, the scalar curvature of a Riemannian manifold $(V,g)$ is denoted by $\Sc_g(V) \colon V \to \R$.

\begin{Proposition} \label{prop:bandwidth} Let $\F$ be a field, let $2 \leq n \leq 7$ and let $0 \leq \ell < n$.
Also assume that one of the following conditions holds:
\begin{enumerate}[\myicon]\itemsep=0pt
 \item $\charac(\F) = 0$,
 \item $\charac(\F) = 2$ and $\ell \leq 2$,
 \item $\charac(\F) > 2$ and $\ell \leq 3$.
\end{enumerate}

Let $(B, g)$ be a compact oriented Riemannian $n$-manifold with boundary satisfying $\Sc_g(B) \geq n(n-1)$.
Let
$
 f \colon B \to [-1, +1] \times {\TTT}^{n- 1 -\ell}
$
be a smooth map satisfying $f(\partial B) \subset \{ -1, +1\} \times {\TTT}^{n- 1-\ell}$ such that there exist classes $c_1, \ldots, c_\ell \in H^{1}(B; \F)$ and a regular value $s \in [-1, +1] \times {\TTT}^{n- 1 -\ell}$ of~$f$ with\footnote{As $f(\partial B) \subset \{-1, +1\} \times {\TTT}^{n-1-\ell}$, the preimage $f^{-1}(s) \subset B$ is a closed smooth oriented submanifold, even if~${s \in \{ -1, +1\} \times \TTT^{n-1-\ell}}$.}
\[
\big\langle c_1 \cup \cdots \cup c_\ell, \big[f^{-1}(s) \big]_\F \big\rangle \neq 0 \in \F .
\]
Put $\partial_{\pm} B:= \partial B \cap f^{-1}\big( \{ \pm 1\} \times {\TTT}^{n- 1-\ell} \big) \subset \partial B$.
Then we have \smash{${\rm dist}_g(\partial_- B, \partial_+ B) \leq \frac{2\pi}{n}$}.
\end{Proposition}

\begin{Remark} The assumption $n \leq 7$ in Proposition \ref{prop:bandwidth} and Theorem \ref{thm:tornexist} can be weakened, see Remark \ref{rem:higherdim_mantra}.
\end{Remark}

\begin{proof}[Proof of Proposition \ref{prop:bandwidth}]
Without loss of generality, we can assume that $B$ is connected.
Suppose ${\rm dist}_g(\partial_- B , \partial_+B) > \frac{2\pi}{n}$.
We will show that this leads to a contradiction.


Up to homotopy relative to $\partial B$, $f$ respects some collar structures near $\partial_+ B \subset B$ and near $\{+1\} \times {\TTT}^{n-1-\ell} \subset [-1 , +1] \times {\TTT}^{n-1-\ell}$.
Let $s \in \{ +1\} \times {\TTT}^{n-1-\ell}$ be a regular value of $f|_{\partial_+ B} \colon \partial_+ B \to \{ +1\} \times {\TTT}^{n-1-\ell}$.
Then $s$ is also a regular value of $f$ and $f^{-1}(s) \subset \partial_+ B$.
So, by assumption and setting $\vartheta := c_1 \cup \cdots \cup c_{\ell} \in H^{\ell} (B; \F)$, we obtain
\begin{equation} \label{regvalbound}
 \big\langle \vartheta , \big[ ( f|_{\partial_+ B})^{-1} (s)\big]_\F \big\rangle \neq 0 \in \F.
\end{equation}
Let $\pi \colon [-1, +1] \times {\TTT}^{n-1-\ell} \to {\TTT}^{n-1-\ell}$ be the projection and let
$
 \omega^{n-1-\ell} \in H^{n-1-\ell}\big( {\TTT}^{n-1-\ell} ; \Z\big)
$
 be the cohomological fundamental class.
By \eqref{regvalbound} and Poincar\'e duality, we have
\begin{equation} \label{boundary}
 \big\langle \big( (\pi \circ f)^* \big(\omega^{n-1-\ell}\big) \big) \cup \vartheta , [ \partial_+ B]_\F \big\rangle \neq 0 \in \F.
\end{equation}

The triple $(B, \partial_- B, \partial_+ B)$ is a {\em band} in the sense of \cite[Definition 1.1]{Raede22}.
We claim that there exists a closed oriented smooth hypersurface $\Sigma^{n-1} \subset {\rm int}(B)$ which separates $\partial_- B$ and~$\partial_{+}B$ and admits a metric $g^{n-1}$ of positive scalar curvature.
Indeed, if this is not the case, then the $\frac{2\pi}{n}$-in\-equality \cite[Theorem 2.8]{Raede22}, see also \cite[Sections 3.6, 3.7 and 5.2]{Gr21}, implies that ${\rm dist}_g(\partial_- B , \partial_+B) \leq \frac{2\pi}{n}$, which contradicts our assumption.
Hence, there exists a hypersurface $\Sigma^{n-1}$ with the stated properties, and in particular it follows that $n \geq 3$.

Possibly after removing some components, the hypersurface $\Sigma^{n-1}$ is homologous to $\partial_+ B$ in~$B$, that is,
\begin{equation} \label{homologies}
 \big[\Sigma^{n-1}\big] = [\partial_+ B] \in H_{n-1}(B; \Z).
\end{equation}
Consider the composition
\[
 \Psi^{n-1} \colon\ \Sigma^{n-1} \stackrel{f|_{\Sigma^{n-1}}}{\longrightarrow} [ -1, +1] \times {\TTT}^{n-1-\ell} \stackrel{\pi}{\longrightarrow} {\TTT}^{n -1-\ell}.
\]
Furthermore, let $ \vartheta_{\Sigma^{n-1}} \in H^{\ell}(\Sigma^{n-1}; \F)$ be the restriction of $\vartheta \in H^{\ell}(B; \F)$.

By \eqref{boundary} and \eqref{homologies}, we have
\[
 \big\langle \big( \Psi^{n-1}\big)^* \big(\omega^{n-1-\ell}\big) \cup \vartheta_{\Sigma^{n-1}} , \big[\Sigma^{n-1}\big]_\F \big\rangle \neq 0 \in \F.
\]
Let $\ell' := \max\{ \ell, 2\}$.
Now we do a {\em Schoen--Yau--Schick descent} to construct closed oriented smooth Riemannian $i$-manifolds $\big(\Sigma^i, g^i\big)$, $\ell' \leq i \leq n-2$, satisfying $\Sc_{g^i} \big( \Sigma^i\big)> 0$ and
\[
 \Sigma^{\ell'} \subset \Sigma^{\ell'+1} \subset \dots \subset \Sigma^{n-2} \subset \Sigma^{n-1} ,
\]
together with maps $ \Psi^i \colon \Sigma^{i} \to {\TTT}^{i - \ell}$ so that, for $\ell' \leq i \leq n-2$,
\begin{enumerate}[(i)]\itemsep=0pt
 \item \label{dual} $\Sigma^{i} \subset \big(\Sigma^{i+1}, g^{i+1}\big)$ is a closed oriented stable minimal smooth hypersurface representing the dual of the cohomology class in $H^1\big(\Sigma^{i+1}; \Z\big)$ represented by the composition
 \[
 \Sigma^{i+1} \stackrel{\Psi^{i+1}}{\longrightarrow} {\TTT}^{i-\ell+1}\longrightarrow \SSS^1,
 \]
 where the second map is the projection onto the second factor in ${\TTT}^{i-\ell+1} = {\TTT}^{i-\ell} \times \SSS^1$.
 \item \label{restpsi} $\Psi^i$ is given by the composition
 \smash{$
 \Psi^{i} \colon\ \Sigma^{i} \stackrel{\Psi^{i+1}}{\longrightarrow} {\TTT}^{i-\ell+1} \longrightarrow {\TTT}^{i-\ell}
 $}
 where the second map is the projection onto the first factor in ${\TTT}^{i-\ell+1} = {\TTT}^{i-\ell} \times \SSS^1$.
 \item The metric $g^i$ is conformal to the metric on $\Sigma^i$ induced by $g^{i+1}$.
 \end{enumerate}
 For details of this construction, see \cite[Theorem 5.1]{S-Y1979} for $i =2$, the proof of \cite[Theorem 1]{SY} for~${i \geq 3}$, and the discussion in \cite{Sch98}.


 Let $\omega^{i-\ell} \in H^{i - \ell}\big( {\TTT}^{i-\ell}; \Z\big)$ be the cohomological fundamental class and let $\vartheta_{\Sigma^{i}}$ be the restriction of $\vartheta_{\Sigma^{i+1}}$ to $\Sigma^i$.
By induction, it follows from \ref{dual} and \ref{restpsi} that
\[
 \big\langle \big(\Psi^{\ell'}\big)^* \big(\omega^{\ell'-\ell}\big) \cup \vartheta_{\Sigma^{\ell'}} , \big[\Sigma^{\ell'}\big] \big\rangle \neq 0 \in \F.
 \]
 In particular, since $\big(\Psi^{\ell'}\big)^* \big(\omega^{\ell'-\ell}\big)$ and \smash{$\vartheta_{\Sigma^{\ell'}}$} are products of classes in \smash{$H^1\big( \Sigma^{\ell'} ; \F\big)$}, we see that \smash{$\big(\Sigma^{\ell'}, g^{\ell'}\big)$} is an $\ell'$-dimensional oriented smooth $\F$-toral \yes-manifold.
After we have passed to an appropriate connected component, we can assume that \smash{$\Sigma^{\ell'}$} is connected.

If $\ell' = 2$, this leads to a contradiction, see Remark \ref{rem:mixed}\,\ref{leqzwei}, so we can assume $\ell' = \ell \geq 3$.
If~${\charac(\F) = 0}$, this implies that $\Sigma^{\ell}$ is overtorical (see Remark \ref{rem:bemerkung}\,\ref{overtorical}), and we can do the Schoen--Yau--Schick descent again to show that $\Sigma^{\ell}$ is \no, contradiction.
If $\charac(\F)$ is finite, we have~${\ell = 3}$ and hence $\charac(\F) > 2$ by assumption.
So we get a contradiction to Proposition~\ref{prop:toralthree}.\looseness=1
\end{proof}

We obtain the following \yes-obstruction.

\begin{Theorem} \label{thm:tornexist} Let $\F$ be a field and $0 \leq \ell \leq n \leq 7$.
Let $V$ be a closed orientable tor-$\ell/\mathbb F$-enlargeable $n$-manifold and assume that one of the following conditions holds:
\begin{enumerate}[\myicon]\itemsep=0pt
 \item $\charac(\F) = 0$,
 \item $\charac(\F) = 2$ and $\ell \leq 2$,
 \item $\charac(\F) > 2$ and $\ell \leq 3$.
\end{enumerate}
Then $V$ is \no.
\end{Theorem}



\begin{proof} Suppose there is a counterexample $(V,g)$ with $\Sc_g(V) > 0$ of dimension $2 \leq n \leq 7$.
Choose an orientation of $V$.
By a scaling of $g$ we can assume $\Sc_g(V) \geq n(n-1)$.

If $\ell = n$, then a finite Riemannian covering $\bar V$ of $V$ is $\mathbb{F}$-toral.
If $\charac(\F) = 0$ then there exists a~map of non-zero degree $\bar V \to {\TTT}^n$, see Example \ref{rem:bemerkung}\,\ref{overtorical}.
Applying a Schoen--Yau--Schick descent as in the proof of Proposition \ref{prop:bandwidth}, this implies that $\bar V$ is \no.
If $\charac(\F) = 2$ and~${n = \ell = 2}$, then~$\bar V$ is \no~by Remark \ref{rem:mixed}\,\ref{leqzwei}.
If $\charac(\F) > 2$ and $n = \ell \leq 3$, then $\bar V$ is \no, by Proposition~\ref{prop:toralthree} for $n =3$, or Remark \ref{rem:mixed}\,\ref{leqzwei} for $n = 2$.
In summary, for $\ell = n$ we see that $\bar V$ and hence $V$ are~\no.

So we can assume that $\ell < n$.
Choose a smooth embedding
\[
 \psi \colon\ [-1,+1] \times \TTT^{n-1-\ell} \hookrightarrow \SSS^{n-\ell} \setminus \{\infty \},
\]
where $[-1,+1] \times \TTT^{n-1-\ell}$ carries the product of the standard metrics.
For $\delta \in (0,1]$, we obtain the smooth submanifold $A_\delta := \psi\big( [ - \delta , + \delta] \times \TTT^{n-1-\ell} \big) \subset \SSS^{n-\ell}$ with boundary
\[
 \partial A_\delta = \partial_- A_\delta \cup \partial_+ A_\delta , \qquad \partial_{\pm} A_\delta = \psi\big(\{ \pm \delta\} \times \TTT^{n-1-\ell}\big) .
\]
We equip $A_\delta$ with the Riemannian metric induced from $\SSS^{n-\ell}$.
Since $A_1 \subset \SSS^{n-\ell}$ is compact, the inverse $\psi^{-1} \colon A_1 \to [-1,+1] \times \TTT^{n-1-\ell}$ is $\Lambda$-Lipschitz for some $\Lambda > 0$.
Hence, for each $\delta \in (0,1]$, we have
\begin{equation} \label{estdist}
 2 \delta \leq \Lambda \cdot {\rm dist}_{A_\delta} ( \partial_- A_\delta, \partial_+ A_\delta) .
\end{equation}
By assumption, there exist a connected cover $\bar V \to V$, cohomology classes $c_1, \ldots , c_{\ell} \in H^1\big( \bar V; \F\big)$, an $\eps := \tfrac{n}{2 \pi \Lambda}$-contracting compactly supported map $f \colon \bar V \to \SSS^{n-\ell}$ and a regular value $s \in A_{\nicefrac{1}{2}} \subset \SSS^{n-\ell} \setminus \{ \infty\}$ of $f$ such that
\begin{equation} \label{regval}
 \big\langle c_1 \cup \cdots \cup c_{\ell}, \big[ f^{-1}(s) \big]_\F \big\rangle \neq 0 \in \F .
\end{equation}
Let $\delta \in \big(\tfrac{1}{2} ,1\big)$ be such that $f$ is transverse to $\partial A_\delta \subset \SSS^{n-\ell}$ and let
$
 B := f^{-1} ( A_\delta) \subset \bar V $.
This is a~smooth compact oriented submanifold of codimension $0$ with boundary
$
 \partial B = \partial_- B \cup \partial_+ B$, $ \partial_{\pm} B := f^{-1} (\partial_{\pm} A_\delta) $.
We equip $B$ with the Riemannian metric from $V$.
By our choice of $\varepsilon$ and using \eqref{estdist} together with $\delta > \tfrac{1}{2}$, we have
\begin{equation} \label{breite}
 {\rm dist}_B( \partial_- B, \partial_+ B) \geq \frac{2 \pi \Lambda}{n} \cdot {\rm dist}_{A_\delta} (\partial_- A_\delta, \partial_+ A_\delta)
 \geq \frac{2 \pi \Lambda}{n} \cdot \Lambda^{-1} \cdot 2\delta > \frac{2\pi}{n} .
\end{equation}
Let $c'_1, \ldots , c'_\ell \in H^1(B; \F)$ denote the restrictions of $c_1, \ldots , c_\ell \in H^1\big(\bar V; \F\big)$.
We obtain by \eqref{regval}
\[
 \big\langle c'_1 \cup \cdots \cup c'_{\ell}, \big[ f^{-1}(s) \big]_{\F} \big\rangle \neq 0 \in \F.
\]
Using $\Sc_g(B) \geq n(n-1)$ and \eqref{breite}, this contradicts Proposition \ref{prop:bandwidth} since $A_{\delta}$ is diffeomorphic to~${[-1, +1] \times \TTT^{n-1-\ell}}$.
\end{proof}

\begin{Remark} \label{rem:higherdim_mantra}\qquad
 \begin{enumerate}[(a)]\itemsep=0pt
\item \label{mantra1} If $n = 8$, then \cite{Smale} implies that the $\mu$-bubble construction underlying the proof of \cite[Theorem 2.8]{Raede22} can be performed after a small $C^{\infty}$-perturbation of $g$.
In view of the recent results in \cite{CMS}, this can also be achieved for $n = 9, 10$.
Higher dimensions $n \geq 11$ may possibly be covered by the results in \cite{Loh2018}.
Thus the condition $n \leq 7$ in Theorem \ref{thm:tornexist} can be weakened accordingly.
\item The real projective spaces $\R P^n$ are $\F_2$-toral for all $n$, are \yes~for $n \geq 2$, and orientable for all odd $n \geq 3$.
Thus, if $\charac(\F) = 2$, then neither the orientability of $V$ nor the assumption~${\ell \leq 2}$ in Theorem \ref{thm:tornexist} can be dropped.
\end{enumerate}
\end{Remark}

\begin{Example} \label{ex:counterexample} Let $4 \leq n \leq 7$, let $0 \leq \ell \leq n $ and let the orientable manifold $W$ be obtained from~${{\TTT}^n = {\TTT}^{\ell} \times {\TTT}^{n-\ell}}$ by killing the $2$-multiples of generators of $\pi_1\big({\TTT}^\ell\big) = \Z^\ell$ by surgery along pairwise disjoint embedded circles in ${\TTT}^n$.
Let $V$ be the connected sum of $W$ and $\mathbb{C} P^{2} \times \SSS^{n-4}$.

Then $V$ is an orientable tor-$\ell / \F_{\!2}$-enlargeable manifold (compare Example \ref{ex:torsion}) which is~\yes, if and only if $\ell \geq 3$.

The \no~property for $\ell \leq 2$ follows from Theorem \ref{thm:tornexist} while the \yes~property for $3 \leq \ell \leq n$ is shown as follows:
We have $\pi_1(V) = (\Z/2)^{\ell} \times \Z^{n-\ell}$ and the classifying map $\varphi \colon V \to \sB \pi_1(V)$ sends the integral fundamental class of $V$ to the cross product of the toral homology classes in~${H_{\ell}\big((\Z/2)^{\ell} ; \Z\big)}$ and $H_{n-\ell}\big(\Z^{n-\ell}; \Z\big)$.
Since $\ell \geq 3$, by \cite[Theorem A]{Jo2004} the toral homology class in~${H_{\ell}\big((\Z/2)^{\ell} ; \Z\big)}$ can be represented by a closed oriented \yes~$\ell$-manifold, and so the same holds for $\varphi_*([V]) \in H_n(\pi_1(V) ; \Z)$.
Using the fact that the universal cover of $V$ is non-spin, Theorem~\ref{prop:nonspin} implies that $V$ is \yes.
\end{Example}


The following fundamental conjecture remains open.

\begin{Conjecture} \label{main_conj}
Let $\F$ be a field of odd characteristic, let $n \geq 4$ and let $4 \leq \ell \leq n$.
Then each closed tor-$\ell/\mathbb F$-enlargeable $n$-manifold is \no.
\end{Conjecture}

Recall that Corollary \ref{cor:asphthree} is specific to $3$-dimensional manifolds, see Remark \ref{rem:mixed}\,\ref{notinge4}, and thus our proof of Theorem \ref{thm:tornexist} does not generalize to $\ell \geq 4$ if $\charac(\F) > 2$.

Table \ref{tab:conditions} collects our knowledge about the existence of closed orientable tor-$\ell / \F$-enlargeable \yes-manifolds of dimension $n \leq 7$ (or more general $n$ as in Remark \ref{rem:higherdim_mantra}\,\ref{mantra1}).

\begin{table}[h]\renewcommand{\arraystretch}{1.2}
\centering
\begin{tabular}{|c|c|c|c|}
\hline
 & $\charac(\F)$ = 0 & $\charac(\F) = 2$ & $\charac(\F) > 2$ \\
\hline
$\ell \leq 2$ & no & no & no \\
\hline
$\ell = 3$ & no & yes & no \\
\hline
$\ell > 3$ & no & yes & \textcolor{red}{open} \\
\hline
\end{tabular}

\caption{\label{tab:conditions}Existence of \yes-manifolds.}
\end{table}

\section[PSC-products]{\yes-products of \no-manifolds}

Given a closed connected oriented manifold $V$, we denote by
\[
 \Phi_{V} = \varphi_*([V]) \in \Ho_*( \pi_1(V);\Z)
\]
 the image of the homological fundamental class of $V$ under the classifying map $\varphi \colon V \to \sB \pi_1(V)$.


As an application of our discussion we obtain many \no-manifolds whose Cartesian products are \yes.

\begin{Theorem}\label{thm:plenty}
Let $V_1$, $V_2$ be closed oriented connected manifolds such that
\begin{enumerate}[\myicon]\itemsep=0pt
 \item $\dim (V_1 \times V_2) \geq 5$;
 \item the universal cover of $V_1 \times V_2$ is non-spin;
 \item $\Phi_{V_1} \in \Ho_*(\pi_1(V_1) ; \Z)$ and $\Phi_{V_2} \in H_*(\pi_1(V_2); \Z)$ are of finite, coprime order.
\end{enumerate}
Then $V_1 \times V_2$ is \yes.
\end{Theorem}

\begin{proof} The class $\Phi_{V_1 \times V_2} \in \Ho_*(\sB ( \pi_1(V_1) \times \pi_1(V_2)); \Z)$ is the image of $\Phi_{V_1} \otimes \Phi_{V_2}$ under the cross-product homomorphism
\[
 \Ho_*( \pi_1(V_1); \Z) \otimes \Ho_* ( \pi_1(V_2); \Z) \to \Ho_* ( \pi_1(V_1) \times \pi_1(V_2); \Z) .
\]
Hence, this class is equal to $0$, since $\Phi_{V_1}$ and $\Phi_{V_2}$ are elements of coprime order in the left-hand tensor factors.
Now the claim follows from Theorem \ref{prop:nonspin}.
\end{proof}

\begin{Example} \label{interesting}
Let $p_1$, $p_2$ be different odd primes, let $4 \leq n \leq 7$, let $1 \leq \ell \leq 3$ and let~$V_1$,~$V_2 $ be tor-$\ell /\F_{p_i}$-enlargeable manifolds of dimension $n$ obtained from ${\TTT}^{n} = {\TTT}^{\ell} \times {\TTT}^{n-\ell}$ by killing the $p_i$-multiples of the $\ell$ generators in $\pi_1\big({\TTT}^\ell\big) =\Z^{\ell}$ by surgery, see Example \ref{ex:torsion}.
Replace $V_1$ by~${V_1 \sharp \big( \mathbb{C}P^2 \times \SSS^{n-4} \big)}$ in order to make sure that the universal cover of $V_1 \times V_2$ is non-spin.
Then
\begin{enumerate}[\myicon]\itemsep=0pt
 \item $V_1$ and $V_2$ are \no~by Theorem \ref{thm:tornexist}.
 \item $V_1 \times V_2$ is \yes~by Theorem \ref{thm:plenty} and the final statement in Example \ref{ex:torsion}.
\end{enumerate}
\end{Example}

\begin{Remark}\quad
\begin{enumerate}[(a)]\itemsep=0pt
\item Examples \ref{interesting} elaborates on the $\mathcal{SYS} \times \mathcal{SYS}$-remark from \cite{Gr2018}.
They are the first examples of their kind in dimensions greater than $4$ and with non-spin universal covers.
\item Complex surfaces of odd degree $\geq 5$ in $\mathbb{C}P^3$ are \no~by Seiberg--Witten theory, while their products are \yes, as are all simply connected non-spin $n$-manifolds for $n\geq 5$, see~\cite{GL}.
\item A referee pointed out the following class of spin examples.
Recall the graded coefficient ring of connective real $K$-homology,
\[
 {\rm ko}_* \cong \Z[ \eta, \omega, \mu] / \big( 2 \eta, \eta^3, \omega \eta, \omega^2 - 4 \mu\big), \qquad |\eta| = 1,\qquad |\omega| = 4,\qquad |\mu|=8.
\]
By \cite[Theorem A]{Stolz92}, a simply connected closed spin manifold of dimension $n \geq 5$ is~\yes~if and only if the Hitchin index \cite{Hitchin} vanishes, $\alpha(M) = 0 \in {\rm ko}_n$.
Since ${\rm ko}_*$ has zero divisors, one obtains simply connected closed spin \no~manifolds $V_1$ and $V_2$ whose products are~\yes.

For example, it is well known, see \cite[p.\ 44]{Hitchin}, that for $n \equiv 1, 2 \mod(8)$, there are exotic spheres $\Sigma^n$ with $\alpha(\Sigma^n) \neq 0 \in {\rm ko}_n \cong \Z/2$.
Let $K^4$ be the K3 surface.
It is simply connected and satisfies $\alpha(K^4) = \omega \in {\rm ko}_4 \cong \Z$.
Since $\eta \cdot \omega = 0$ and the Hitchin index is multiplicative, the products $\Sigma^n \times K^4$ are \yes, whereas $\Sigma^n$ and $K^4$ are \no.
\end{enumerate}
\end{Remark}

\section{Four conjectures} \label{fourconjectures}

\begin{Conjecture}[short/long neck torality inequality\footnote{This is motivated by Cecchini's long neck inequality \cite{Ce2020}.}]
\label{shortneck}
Let $0 \leq \ell < n$ be natural numbers, let~$(V,g)$ be a compact oriented Riemannian $n$-manifold with boundary and let
$f \colon  V\to \SSS^{n-\ell}$
be a smooth {\it area decreasing} map\footnote{That is ${\rm area}(f(\Sigma)) \leq {\rm area}(\Sigma)$ for all smooth surfaces $\Sigma\subset X$.} which is constant on the boundary $\partial V\subset V $, say $f(\partial V)=\{ \infty \}\subset \SSS^{n-\ell}$.
Suppose $\charac(\F) \neq 2$ and there are classes $c_1, \ldots, c_\ell \in H^1(V; \F)$ such that for some regular value $s \in \SSS^{n-\ell} \setminus \{ \infty\}$ of $f$ we get
\[
 \bigl\langle c_1 \cup \cdots \cup c_\ell, \bigl[ f^{-1}(s)\bigr] \bigr\rangle \neq 0 \in \F .
\]
Setting $k := n - \ell$, suppose that for all $x \in \operatorname{supp}( {\rm d}f) = \{ x \in V \mid {\rm d}f(x) \neq 0\}$, we have
\[
 \Sc_g(V, x) \geq k(k-1) .
\]
If the scalar curvature of $V$ is bounded from below by $\sigma > 0$, i.e.,
\[
 \Sc_g(V,x)\geq \sigma>0 \qquad \textrm{for } x\in V,
 \]
then the distance from the support of ${\rm d}f$ to the boundary of $V$ satisfies
\[
 {\rm dist}_g ( \operatorname{supp}({\rm d}f), \partial V) < {\rm const} \cdot \sqrt{\tfrac {1}{\sigma}},
 \]
 possibly for \smash{${\rm const}=\pi \sqrt{\frac{n-1}{n}}$} where this constant is optimal.
 \end{Conjecture}

 \begin{Remark} There is a significant mismatch between distance and area domination with positive scalar curvature.
Namely, given an orientable $n$-manifold $V=(V,g)$ with $\Sc_g(V) \geq \sigma>0$, then, for all $\varepsilon >0$, there exists a metric $g_{\varepsilon}$ on $V$, such that
\begin{enumerate}[\myicon]\itemsep=0pt
\item $ V_{\varepsilon}=(V,g_{\varepsilon})$ admits an {\it area decreasing} diffeomorphism $V\to V_{\varepsilon }$;
\item no Riemannian manifold $(W,h)$ which $1$-Lipschitz dominates $V_\eps$, i.e., which admits a proper $1$-Lipschitz map of non-zero degree $W \to V_\eps$, has $\Sc_h \geq \varepsilon$.
\end{enumerate}
The simplest example is where $V=\SSS^2\times \TTT^{n-2}$ and $V_\varepsilon =\SSS^2_{\varepsilon, D}\times \TTT_\varepsilon^{n-2}$, where $\SSS^2_{\varepsilon, D}$ is the (smoothed) boundary of the cylinder with base of radius $\varepsilon$ and height \smash{$D=10\sqrt{ \tfrac {1}{\varepsilon}}$}, i.e.,
\[
 \SSS^2_{\varepsilon, D}=\partial \big( B^2(\varepsilon)\times [0,D] \big),
 \]
$\TTT^{n-2} = \SSS^1 \times \cdots \times \SSS^1$ is the $(n-2)$-torus with the standard metric and $\TTT_\varepsilon^{n-2}$ is the product of circles with radius $\varepsilon$.
Here, the existence of an area-decreasing diffeomorphism $V\to V_{\varepsilon, }$ (for~${\varepsilon \ll 0.01}$) is elementary, while the non-existence of $1$-Lipschitz domination with $Sc \geq \varepsilon$ follows from the~$\frac {2\pi}{n}$-inequality as $10 > 2\pi$, see \cite[Section 3.6]{Gr21} and \cite[Theorem 2.8 and Proposition~2.21]{Raede22}.\footnote{The case of $n \leq 7$ is classical, while the case $8 \leq n \leq 10$ uses the methods of \cite{Smale} and \cite{CMS}, respectively.
For $n \geq 11$, see \cite[Theorem 4.6]{S-Y2017} and \cite{Loh2018}.
If $V_\eps$ is spin, then the full $\frac {2\pi}{n}$-inequality is proved in \cite[Theorem 1.4]{Ze2019}, \cite[Corollary 1.5]{Ze2020}, \cite[Theorem D]{Ce2020} and \cite[Theorem 1.3]{G-X-Y2020}.}

Now, for an arbitrary $(V,g)$, let $\TTT^{n-1}\subset V $ be a smoothly embedded torus and let us modify the metric $g$ in a small neighborhood $\TTT^{n-1}\times [-\delta, +\delta]\subset V$ of the torus by $\varepsilon'$-shrinking $g$ along the torus and \smash{$\sqrt \frac {1}{\varepsilon'}$}-stretching $g$ normally to the torus.
For sufficiently small $\eps' = \eps'(g, \eps) > 0$, this decreases the areas of the surfaces in $V$, while the $\frac {2\pi}{n}$-inequality prevents $1$-Lipschitz domination of $V_\varepsilon$ with $\Sc \geq \varepsilon$.
\end{Remark}

\begin{Conjecture}[tor-non-enlargeability of foliations with $\Sc>0$] \label{foliations}
 Let $V$ be a closed $n$-manifold and $\mathscr F$ be a smooth foliation on $V$ equipped with a leafwise Riemannian metric $g_{\mathscr F}$, i.e., $g_{\mathscr F}$ is a smooth bundle metric on the integrable subbundle $T \mathscr F \subset TV$ tangent to the leaves of $\mathscr F$.

Let $\Sc(\mathscr F) \colon V \to \R$ be defined as $\Sc(\mathscr F, x) = \Sc_{g_{\mathscr F}}({\mathscr F}_x , x)$ where ${\mathscr F}_x$ is the leaf passing through $x \in V$ equipped with the Riemannian metric induced by $g_{\mathscr F}$.


If $V$ is tor-$\ell / \mathbb F$-enlargeable for some non-negative integer $\ell$ and some field $\F$ of $\charac(\F) \neq 2$, then $\Sc(\mathscr F) \ngtr 0$.
 \end{Conjecture}

For foliations of codimension $1$, we have the following partial answer.

\begin{Proposition} Let $\F$ be a field, let $1 \leq n \leq 6$ and let $0 \leq \ell \leq n$.
Also assume that one of the following conditions holds:
\begin{enumerate}[\myicon]\itemsep=0pt
 \item $\charac(\F) = 0$,
 \item $\charac(\F) = 2$ and $\ell \leq 2$,
 \item $\charac(\F) > 2$ and $\ell \leq 3$.
\end{enumerate}
Let $V$ be a closed orientable tor-$\ell/\F$-overtorical\footnote{See Example \ref{ex:toral}\,\ref{ex:overtorical}.} $n$-manifold.

Then $V$ does not admit a smooth orientable codimension-$1$ foliation $\mathscr F=(\mathscr F, g_{\mathscr F})$ together with a leafwise Riemannian metric such that $ \Sc (\mathscr F)>0$.
\end{Proposition}

\begin{proof}[Sketch of proof] Suppose $V$ carries a smooth leafwise Riemannian metric $g_{\!\mathscr F}$ with ${\Sc_{g_{\!\mathscr F}} (\mathscr F)\!\! > \!\!0}$.
After scaling $g_{\mathscr F}$, we can assume that $\Sc_{g_{\mathscr F}} (\mathscr F) > (n+1)n$.
Let $V^* \to V $ be the (one-dimensional for $\mathrm{codim} (\mathscr F)=1$) {\it Connes bundle} of $\mathscr F$, which in this case is equivalent to the trivial $\R$-bundle~${V\times \mathbb R \to V}$, see \cite[Sections $1\frac{7}{8}$ and $9\frac {2}{3}$]{Gr1996} and \cite[Section 6.5.2]{Gr21}.
Recall \cite[p.\ 9]{Gr1996} that the total space $V^*$ of this bundle has an induced codimension-$2$ foliation $\mathscr F^*$ whose leaves correspond to leaves in $\mathscr F$, and that $V^*$ has a Riemannian metric $\tilde g$ such that
\begin{enumerate}[\myicon]\itemsep=0pt
 \item on each leaf of $\mathscr F^*$, the metric induced by $\tilde g$ coincides with the metric induced by $g_{\mathscr F}$ on the corresponding leaf of $\mathscr F$,
 \item the restriction of $\tilde g$ to $(T \mathscr F^*)^{\perp} \subset TV^*$ is {\em transversally Riemannian}, i.e., it is a bundle metric which is parallel with respect to the Bott connection on $(T \mathscr F^*)^{\perp} \to V^*$ along $\mathscr F^*$.
\end{enumerate}

For $\lambda > 0$, we define a new Riemannian metric $\tilde g_{\eps}$ on $V^* \cong V \times \R$ as the bundle metric on
\[
{TV^* = T\mathscr F^* \oplus (T\mathscr F^*)^{\perp}}
\]
 equal to
\[
 \tilde g_\lambda = \tilde g|_{T\mathscr F^*} + \lambda^2 \tilde g|_{(T\mathscr F^*)^{\perp}}.
\]
The discussion in \cite[p.\ 315]{Gr21} shows that there exist $\lambda \gg 0$ such that
\begin{enumerate}[\myicon]\itemsep=0pt
\item $\Sc_{\tilde g_{\lambda}} (V \times [-1, +1]) \geq (n+1)n$,
\item ${\rm dist}_{\tilde g_\lambda}( V \times \{-1\}, V\times \{+1\}) > \frac{2\pi}{n+1}$.
\end{enumerate}
Since $V$ is tor-$\ell/\F$-overtorical, this contradicts Proposition \ref{prop:bandwidth}.
\end{proof}

The above suggests the following.

\begin{Conjecture}[domination of foliations] If a closed orientable manifold $V$ admits a smooth foliation $\mathscr F$ with $\Sc(\mathscr F) > 0$, then $V$ can be dominated by an orientable manifold $W$ with $\Sc(W)>0$.\footnote{This is hardly to be expected from {\it leaf-wise} smooth Riemannian foliations $\mathscr F$ which are (only) {\it continuous} transversally to the leaves.
However, the standard examples of such $\mathscr F$ -- the stable/unstable foliations of {\it Anosov systems} -- are easily shown to be \no.}
\end{Conjecture}

We conclude by bringing together Conjectures \ref{shortneck} and \ref{foliations} as follows, compare \cite{S-W-Z2021}.

\begin{Conjecture}[foliated short/long neck torality inequality]
Let $V$ be a compact oriented Riemannian $n$-manifold with boundary and $\mathscr F$ be a smooth $m$-dimensional foliation on $V$.

Let $\charac(\F) \neq 2$ and let $f \colon V \to \SSS^{n-\ell}$ be a smooth map which {\it decreases the areas of smooth surfaces in the leaves of} $\mathscr F$ such that $f(\partial V)=\{ \infty\} \subset \SSS^{n-\ell}$ and such that there are classes~${c_1 , \ldots, c_\ell \in H^1(V ; \F)}$ with
\[
 \big\langle c_1 \cup \cdots \cup c_{\ell}, \big[f^{-1}(s)\big] \big\rangle \neq 0 \in \F
\]
for some regular value $s \in \SSS^{n-\ell} \setminus \{ \infty\}$.

Setting $k := n-\ell$, let
\[
 \Sc(\mathscr F,x)\geq\min \left\{ k(k-1), m(m-1)\right\}
\]
for all $x\in V$ with $x \in \operatorname{supp}( {\rm d}f_{\mathscr F}) = \{ x \in V \mid {\rm d}f(x)|_{T \mathscr F} \neq 0\}$.


If the scalar curvature of $\mathscr F$ is bounded from below by $\sigma > 0$, i.e.,
\[
 \Sc(\mathscr F,x)\geq \sigma>0 \qquad\text{for}\quad x\in V,
 \]
then {\it the leaf-wise distance} from the support of ${\rm d}f$ on $T\mathscr F$ to the boundary of $V$ satisfies
\[
 {\rm dist}_{\mathscr F}(\operatorname{supp}({\rm d}f_{\mathscr F}), \partial V)=\inf_{x\in \operatorname{supp} {({\rm d}f_{\mathscr F})}} {\rm dist}_{ \mathscr F_x} (x, \mathscr F _x\cap \partial X) < {\rm const} \sqrt{ \tfrac{1}{\sigma}},
\]
where $\mathscr F _x\subset V$ is the leaf of $\mathscr F$ through $x$.
\end{Conjecture}

\appendix


 \section[PSC-metrics on manifolds with non-spin universal covers]{\yes-manifolds with non-spin universal covers}

In \cite{GL}, it was proved that closed simply connected non-spin manifolds of dimension $\geq 5$ are~\yes.
Here we prove the corresponding result for non-simply connected manifolds.


\begin{Theorem} \label{prop:nonspin} Let $n \geq 5$ and let $V$ be a closed connected oriented smooth $n$-manifold with non-spin universal cover.
Let $\varphi \colon V \to {\sf B} \pi_1(V)$ be the classifying map and assume that there exists a map $f \colon V' \to {\sf B}\pi_1(V)$ where $V'$ is a closed oriented $n$-dimensional \yes-manifold such that
\[
 \varphi_*([V]) = f_{*} ( [ V' ]) \in H_n(\pi_1(V) ; \Z) .
\]
$($This includes the case $\varphi_*([V]) = 0$.$)$

Then $V$ is \yes.
\end{Theorem}


The proof is based on the following generic construction of \yes-manifolds, compare \cite{Stolz92} and \cite{Fuehring}.
Consider the action of $G := \SU(3) \rtimes \Z/2$ on $\C P^2$ where $\SU(3)$ acts on $\C P^2$ by matrix multiplication and $\Z/2$ acts by complex conjugation.
This action is transitive, its stabilizer group at $[0:0:1] \in \C P^2$ is equal to $H = {\rm S}({\rm U}(2) {\rm U}(1)) \rtimes \Z/2 < G$, and it is isometric with respect to the Fubini study metric on $\C P^2$.
Hence, we obtain a fibre bundle
\begin{equation} \label{universal}
 \C P^2 = G/ H \hookrightarrow \sB H \to \sB G
\end{equation}
with structure group $G$.
Let $B$ be a~closed oriented smooth manifold, let $X$ be a CW complex and let $
 f \times \omega \colon\ B \to X \times \sB G
$
be a continuous map.
The map $\omega \colon B \to \sB G$ classifies a smooth fiber bundle $\C P^2 \hookrightarrow M_{\omega} \stackrel{\pi}{\longrightarrow} B$ with structure group $G$, which is unique up to fibre bundle diffeomorphism.
The total space $M_{\omega}$ is a closed oriented smooth manifold and it is \yes, by the O'Neill formula.
Furthermore, the assignment $(f \times \omega \colon B \to X \times BG) \mapsto (f \circ \pi \colon M_\omega \to X)$ respects the oriented bordism relation, and we obtain a group homomorphism which is natural in $X$,
\begin{equation} \label{psi}
 \psi(X)_* \colon\ \Omega^{\SO}_{*-4} (X \times \sB G) \to \Omega^{\SO}_{*}(X).
\end{equation}
Let $\mathscr{T}_*(X) := \im ( \psi(X)_*) \subset \Omega^{\SO}_{*}(X)$ and let $ \Omega^{\SO, +}_*( X) \subset \Omega_*^{\SO}(X)$ be the subgroup consisting of oriented bordism classes $[M \to X]$ where $M$ is \yes.
We obtain
$
 \mathscr{T}_*(X) \subset \Omega^{\SO,+}_*(X)$.
Now consider the {\em homological orientation}, which is natural in $X$,
\begin{equation} \label{homor}
 u(X)_* \colon\ \Omega^{\SO}_*(X) \to H_*(X; \Z), \qquad [f \colon M \to X] \mapsto f_*([M]) \in H_*(X; \Z).
\end{equation}
For dimension reasons, $\pi \colon M_\omega \to B$ sends the fundamental class of $M_\omega$ to zero, hence
$
 \mathscr{T}_*(X) \subset \ker u(X)_*$.
We will next show that $2$-locally, the reverse inclusion holds as well.

For an abelian group $A$, we will use the shorthand $A_{(2)} = A \otimes \Z_{(2)}$ for the $2$-localisation.

 \begin{Proposition} \label{crucial} We have
$
 \mathscr{T}_*(X)_{(2)} = (\ker u(X)_*)_{(2)}$.
 In particular, for all $x \in \ker u(X)_*$ there is an odd $k \in \N$ with $kx \in \mathscr{T}_*(X)$.
\end{Proposition}

\begin{proof}
Only the inclusion ``$\supset$'' remains to be shown.
It is well known, see \cite[Chapter 2]{CF}, that the Atiyah--Hirzebruch spectral sequence for $2$-local oriented bordism
\[
 E^2_{p,q} \cong H_p\big(X; \Omega^{\SO}_q\big)_{(2)} \Longrightarrow \Omega^{\SO}_{n}(X)_{(2)}
 \]
collapses at $E^2$.
There is a decreasing filtration
\[
 \Omega^{\SO}_{n}(X)_{(2)} = \mathscr{F}_{n,0} \supset \cdots \supset \mathscr{F}_{0,n} = H_0(X ; \Omega_n)_{(2)}
\]
with \smash{$E^{2}_{p,q} \cong \mathscr{F}_{p,q} / \mathscr{F}_{p-1,q+1}$}.
Furthermore, the usual $\Omega^{\SO}_*$-module structure on $\Omega^{\SO}_*(X)$ induces maps
$
 \Omega^{\SO}_{\mu} \times \mathscr{F}_{p, q} \mapsto \mathscr{F}_{p, q+\mu}
$
so that the isomorphism $E^2_{p,*} \cong H_p\big(X; \Omega^{\SO}_*\big)$ is $\Omega^{\SO}_*$-linear.

The homological orientation $u(X)_*$ corresponds $2$-locally to the projection
\[
 \Omega^{\SO}_{n}(X)_{(2)} = \mathscr{F}_{n,0} \to \mathscr{F}_{n,0} / \mathscr{F}_{n-1,1} = H_n\big(X; \Omega^{\SO} _{0}\big)_{(2)} = H_n(X; \Z_{(2)}).
\]
Thus, by the collapsing of the Atiyah--Hirzebruch spectral sequence, $u(X)_*$ is $2$-locally surjective.

Let $\Tor_* \subset \Omega^{\SO}_*$ be the torsion subgroup.
The short exact sequence
\[
 0 \to \Tor_* \to \Omega^{\SO}_* \to \Omega^{\SO}_* / \Tor_* \to 0
\]
splits, and hence it induces short exact sequences
\begin{equation} \label{ses}
 0 \to H_p(X; \Tor_q)_{(2)} \hookrightarrow E^2_{p,q} \longrightarrow H_p\big(X; \Omega^{\SO}_q / \Tor_q\big)_{(2)} \to 0 .
\end{equation}
It is well known that $\Omega^{\SO}_*$ contains no odd torsion, so we have $H_p(X; \Tor_q)_{(2)} = H_p(X; \Tor_q)$.

The proof of Proposition \ref{crucial} is complete once we show that every element in $E^2_{p,q}$ with $p + q = n$ and $q > 0$ can be represented by an element in $\mathscr{T}_n(X)_{(2)} \cap \mathscr{F}_{p,q}$.
According to \eqref{ses}, we decompose this claim into the following two assertions.


\begin{Assertion}\label{Assertion1} Let $q > 0$.
Then the image of $\mathscr{T}_n(X)_{(2)} \cap \mathscr{F}_{p,q}$ in $E^2_{p,q}$ is mapped surjectively to $H_p\big(X; \Omega_q^{\SO} / \Tor_q \big)_{(2)}$.
\end{Assertion}


To show this, consider a simple tensor
\[
 h \otimes \xi \in H_p\big(X; \Omega_q^{\SO} / \Tor_q \big)_{(2)} = H_p(X; \Z_{(2)}) \otimes \Omega_q^{\SO} / \Tor_q .
\]
There exists an odd $k \in \N$ such that $k\cdot h \in H_p(X; \Z_{(2)})$ lifts to a bordism class $ [ f \colon M^p \to X] \in \Omega_p^{\SO}(X)$.
Furthermore, by constructing appropriate $\C P^2$-bundles, see \cite[Proposition 8.2]{Fuehring_bordism}, one shows that for $q > 0$, the composition
\[
 \Omega^{\SO}_{q-4} (\sB G) \stackrel{\psi({\rm pt.})_*}{\longrightarrow} \Omega^{\SO}_{q} \to \Omega^{\SO}_q/ \Tor_q
\]
is surjective.
Hence, the class $\xi \in \Omega_q^{\SO} / \Tor_q$ is represented by the total space of a fibre bundle with structure group $G$,
\[
 \C P^2 \hookrightarrow N^q \stackrel{\pi}{\longrightarrow} B^{q-4}.
 \]
Now, in \eqref{ses}, the class in $E^2_{p,q}$ represented by
\[
 [ f \colon M^p \to X] \cdot [N^q] = \big[M^p \times N^q \stackrel{\id \times \pi}{\longrightarrow} M^p \times B^{q-4} \stackrel{(x,y) \mapsto f(x)}{\longrightarrow} X\big] \in \mathscr{T}_n(X)_{(2)} \cap \mathscr{F}_{p,q}
\]
maps to $k ( h \otimes \xi) \in H_p\big(X; \Omega^{\SO}_q / \Tor_q \big)_{(2)}$, finishing the proof of Assertion~\ref{Assertion1}.


\begin{Assertion}\label{Assertion2}
 Each element in $H_p(X; \Tor_q) \subset E^2_{p,q}$ is represented by some class in $\mathscr{T}_n(X)_{(2)} \cap \mathscr{F}_{p,q}$.
\end{Assertion}


To show this, we use a homotopy-theoretic argument inspired by \cite{Stolz_splitting}.
For an abelian group~$A$, let ${\sf H}A$ denote the Eilenberg--MacLane spectrum for~$A$, and let $\mathscr{A}^*=H^*({\sf H}\F_2; \F_2)$ be the Steenrod algebra of stable $\F_2$-cohomology operations.
Recall that for any spectrum ${\sf M}$, the cohomology $H^*({\sf M}; \F_2)$ is an $\mathscr{A}^*$-module.

Let $\MSO$ be the oriented bordism spectrum and let $\MSO_{(2)}$ be its $2$-localisation.
Recall that the Pontrjagin--Thom construction yields a natural isomorphism
\[
 \Omega^{\SO}_*(X) \cong \pi_*(\MSO \wedge \Sigma^{\infty} X_+),
\]
where $\wedge$ denotes the smash product and $\Sigma^{\infty} X_+$ is the suspension spectrum of $X$ with a disjoint base point.
From this perspective, the homological orientation $u(X)_*$ from \eqref{homor} corresponds to a spectrum map
\begin{equation} \label{u}
 u \colon\ \MSO \to \mathsf{H}\Z
\end{equation}
and the map $\psi(X)_*$ from \eqref{psi} is identified with
\[
(T \wedge \id_{\Sigma^{\infty} X_+})_* \colon\ \pi_*\big( \Sigma^4 \MSO \wedge \Sigma^{\infty} \sB G_+ \wedge \Sigma^{\infty} X_+\big) \to \pi_*(\MSO \wedge \Sigma^{\infty}X_+) .
\]
Here, the spectrum map $T \colon \Sigma^4 \MSO \wedge \Sigma^{\infty} \sB G_+ \to \MSO$ is the composition
\[
 \Sigma^4 \MSO \wedge \Sigma^{\infty} \sB G_+ \stackrel{\id \wedge \alpha}{\longrightarrow} \Sigma^4 \MSO \wedge {\rm Th}(-T_v \sB H) \stackrel{\id \wedge \beta}{\longrightarrow} \Sigma^4 \MSO \wedge \Sigma^{-4} \MSO \stackrel{\gamma}{\longrightarrow} \MSO,
 \]
where ${\rm Th}(-T_v \sB H)$ is the Thom spectrum of the stable fiberwise normal bundle of \eqref{universal}, $\alpha \colon \Sigma^{\infty} \sB G_+ \to {\rm Th}(-T_v \sB H)$ is the Pontrjagin--Thom map, $\beta \colon {\rm Th}(-T_v \sB H) \to \Sigma^{-4} \MSO$ is the homological orientation, and $\gamma$ is multiplication in the ring spectrum $\MSO$.
For more information on the construction of $T$, compare \cite{Fuehring_bordism}, in particular, Section~3 and Lemma~6.1.

By \cite[Theorem 5]{Wall60}, also compare \cite[proof of~Theorem 14.1]{CF}, there are splittings of spectra
\begin{gather} \label{2localsplit}
 \MSO_{(2)} \simeq \sf{P} \vee \sf{Q} , \qquad \sf{Q} \simeq \mathsf{Q'} \vee \mathsf{H} \Z_{(2)},
 \end{gather}
where $\sf{P}$ is a sum of positive suspensions of ${\sf H}\F_2$ and $\mathsf{Q'}$ is a sum of $4\ell$-fold suspensions of ${\sf H}\Z_{(2)}$, $\ell > 0$.
The induced projection $\MSO_{(2)} \to \mathsf{H}\Z_{(2)}$ is the $2$-localisation of the map ${u \colon \MSO \to \mathsf{H}\Z}$ appearing in \eqref{u}.
In particular, denoting by \smash{$\widehat \MSO \subset \MSO$} the homotopy fibre of $u$, we have~\smash{$ ( \widehat \MSO)_{(2)} \simeq \sf{P} \vee \sf{Q'}$}.
Furthermore, with respect to \eqref{2localsplit}, we have
\begin{equation} \label{splitorbord}
 H^*(\MSO; \F_2) = H^*(\mathsf{P}; \F_2) \oplus H^*( \mathsf{Q}; \F_2) , \qquad \Omega^{\SO}_*(X)_{(2)} = \mathsf{P}_*(X) \oplus \mathsf{Q}_*(X) .
\end{equation}
For $X=\{\rm pt.\}$, the second equality specializes to \smash{$\big(\Omega^{\SO}_*\big)_{(2)} = \pi_*( \mathsf{P}) \oplus \pi_*(\mathsf{Q})$} with $\pi_*(\mathsf{P}) = \Tor_* \subset \Omega^{\SO}_*$.

Since the composition $u \circ T$ represents an element in $H^{-4}( \MSO \wedge \Sigma^{\infty} \sB G_+; \Z) = 0$, it is null-homotopic, and hence the map $T$ lifts to a map \smash{$\widehat T \colon \Sigma^4 \MSO \wedge \Sigma^{\infty} \sB G_+ \to \widehat \MSO$}.
According to \cite[Theorem 2.2]{Fuehring_bordism},\footnote{This reference works with the dual notion of $\mathscr{A}_*$-comodules where $\mathscr{A}_*$ is the dual $\F_2$-Steenrod algebra.} $\widehat T$ induces a split injection of $\mathscr{A}^*$-modules
\[
 \widehat T^* \colon\ H^*\big( \widehat\MSO; \F_2\big) \to H^{*}\big( \Sigma^4 \MSO \wedge \Sigma^{\infty} \sB G_+; \F_2\big) .
\]
Since \smash{$\big( \widehat \MSO\big)_{(2)} \simeq \sf{P} \vee \sf{Q'}$}, this implies that there exists a graded $\mathscr{A}^*$-module $M^*$ and a direct sum decomposition of $\mathscr{A}^{*}$-modules
\begin{equation} \label{splitting1}
 H^{*}\big(\Sigma^4 \MSO \wedge \Sigma^{\infty} \sB G_+; \F_2\big) \cong H^*(\mathsf{P}; \F_2) \oplus M^*,
\end{equation}
so that the restriction of $T^* \colon H^{*} (\MSO; \F_2) \to H^*\big(\Sigma^4 \MSO \wedge \Sigma^{\infty} \sB G_+; \F_2\big)$ to $H^*(\mathsf{P}; \F_2)$, see~\eqref{splitorbord}, is equal to the inclusion $H^*(\mathsf{P}; \F_2) \hookrightarrow H^{*}\big(\Sigma^4 \MSO \wedge \Sigma^{\infty} \sB G_+; \F_2\big)$.

Since $H^*(\mathsf{P}; \F_2)$ is a free $\mathscr{A}^*$-module, according to \cite[Theorem 2\,(a)]{Mar74}, which applies to the locally finite spectrum $\Sigma^4 \MSO \wedge \Sigma^{\infty} \sB G_+$, the splitting \eqref{splitting1} can be realised on the spectrum level, i.e., there is a spectrum $\mathsf{M}$ with $H^*(\mathsf{M} ;\F_2) \cong M^*$ as $\mathscr{A}^*$-modules and a homotopy equivalence (of unlocalized spectra)
\begin{equation} \label{Margolis1}
\Sigma^4 \MSO \wedge \Sigma^{\infty} \sB G_+ \simeq \mathsf{P} \vee \mathsf{M},
 \end{equation}
which induces the splitting \eqref{splitting1} in $\F_2$-cohomology.
Consider the map
\[
 \Xi \colon\ \mathsf{P} \stackrel{\eqref{Margolis1}}{\hookrightarrow} \Sigma^4 \MSO \wedge \Sigma^{\infty} \sB G_+ \stackrel{T}{\longrightarrow} \MSO \to \MSO_{(2)} \stackrel{\eqref{2localsplit}}{\simeq} \mathsf{P} \vee \mathsf{Q} .
 \]
In $\F_2$-cohomology, the composition of this map with the projection onto $\mathsf{P}$ induces the identity on $H^*(\mathsf{P}; \F_2)$, while the composition with the projection onto $\mathsf{Q}$ induces the $0$-map.
Hence, by \cite[Theorem 3]{Mar74} and using that $\mathsf{P}$ is a sum of suspensions of ${\sf H}\F_2$, the map $\Xi$ is homotopic to the inclusion $\mathsf{P} \to \mathsf{P} \vee \mathsf{Q}$ onto the first summand.

We conclude that the composition
\[
 \mathsf{P}_*(X) \to \pi_*\big( \Sigma^4 \MSO \wedge \Sigma^{\infty} \sB G_+ \wedge \Sigma^{\infty}X_+\big) \stackrel{\psi(X)_*}{\longrightarrow} \pi_*(\MSO \wedge \Sigma^{\infty} X_+) = \mathsf{P}_*(X) \oplus \mathsf{Q}_*(X)
\]
is equal to the inclusion $ \mathsf{P}_*(X) \subset \mathsf{P}_*(X) \oplus \mathsf{Q}_*(X)$.
In particular, in the decomposition \eqref{splitorbord}, we have $\mathsf{P}_*(X) \subset \mathscr{T}_*(X)_{(2)}$.

The splitting $\MSO_{(2)} \simeq \mathsf{P} \vee \mathsf{Q}$ in \eqref{2localsplit} induces a splitting $E^r_{p,q} = E^{\mathsf{P}, r}_{p,q} \oplus E^{\mathsf{Q}, r}_{p,q}$ of the Atiyah--Hirzebruch spectral sequence for $\Omega^{\SO}_*(X)_{(2)}$.
Clearly, the Atiyah--Hirzebruch spectral sequences for $\mathsf{P}$ and $\mathsf{Q}$ collapse at $E^2$, and we obtain an induced decomposition
\[
 H_p\big(X; \Omega^{\SO}_q\big)_{(2)} = E^2_{p,q} = E(\mathsf{P})^{2}_{p,q} \oplus E(\mathsf{Q})^{2}_{p,q} = H_p(X ; \pi_q( \mathsf{P})) \oplus H_p(X ; \pi_q( \mathsf{Q})),
\]
where
\[
 E(\mathsf{P})^{2}_{p,q} = \left( \mathscr{F}_{p,q} \cap \mathsf{P}_n(X)\right) / \left( \mathscr{F}_{p-1,q+1} \cap \mathsf{P}_n(X) \right).
\]
Since $H_p(X ; \pi_q( \mathsf{P})) = H_p(X; \Tor_q) \subset E^2_{p,q}$ and $\mathsf{P}_n(X) \subset \mathscr{T}_n(X)_{(2)}$, this finishes the proof of Assertion~\ref{Assertion2}.
\end{proof}

We obtain the following bordism theoretic description of $H_*(X ; \Z_{(2)})$ whose analogue for spin bordism and ${\rm ko}$-homology was shown in \cite[Theorem B\,(2)]{Stolz_splitting}.

\begin{Corollary} For each CW complex $X$, the orientation $u(X)_*$ induces an isomorphism
\[
 \left( \Omega^{\SO}_*(X) / \mathscr{T}_*(X) \right)_{(2)} \cong H_*(X; \Z)_{(2)} .
\]
\end{Corollary}

\begin{proof}[Proof of Theorem \ref{prop:nonspin}]
By the {\em bordism principle} \cite[Theorem 1.5]{R-S1995}, which follows from the {\em surgery principle} in scalar curvature geometry \cite{GL,SY}, and since the universal cover of $V$ is non-spin, it would suffice to show that the oriented bordism class
\[
 [ \varphi \colon V \to \sB \pi_1(V)] \in \Omega^{\SO}_*( \sB \pi_1(V))
\]
lies in $\Omega^{\SO, +}_*( \sB \pi_1(V))$.
We will achieve this goal for an odd multiple of $[ \varphi \colon V \to {\sf B} \pi_1(V)]$, while the rest of the proof relies on a version of the bordism principle for manifolds with Baas--Sullivan singularities worked out in \cite{Han20}.

By assumption, the class
\[
 \beta := [ \varphi \colon V \to \sB \pi_1(V) ] - [f \colon V' \to \sB \pi_1(V) ] \in \Omega_*( \sB \pi_1(V))
\]
is in the kernel of $u(\sB \pi_1 (V))_* \colon \Omega_*(\sB \pi_1(V)) \to H_*(\pi_1(V) ; \Z)$.
By Proposition \ref{crucial}, there is an odd $k \in \N$ with $k \beta \in \mathscr{T}_*(\sB \pi_1(V)) \subset \Omega^{\SO, +} ( \sB \pi_1(V)) $.
Write $k \beta = [ M \to \sB \pi_1(V)]$ where $M$ is~\yes.
Since $[f \colon V' \to \sB \pi_1(V) ] \in \Omega^{\SO,+}_*( \sB \pi_1(V))$ by assumption, we obtain
\[
 k\cdot [ \varphi \colon V \to \sB \pi_1(V) ] \in \Omega^{\SO, +}_*( \sB \pi_1(V)) .
\]

We now use the bordism principle for \yes-manifolds based on {\em manifolds with Baas--Sullivan singularities} \cite[Proposition 3.11]{Han20}.
Let $\Omega^{\SO, \mathscr{Q}}_*(-)$ denote oriented bordism with Baas--Sullivan singularities in $\mathscr{Q}$, where $\mathscr{Q} = (Q_0 = *, Q_1, Q_2, \ldots)$ is a family of closed oriented manifolds~$Q_i$ of dimension $4i$ which are equipped with fixed positive scalar curvature metrics and such that~${\Omega^{\SO}_* / \Tor = \Z [[Q_1], [Q_2], \ldots]}$.
Note that $\mathscr{Q}$ is a positive family of singularity types in the sense of \cite[Definition 3.2]{Han20}, also compare the remarks before \cite[Definition 3.12]{Han20}.

Recall from \cite{baas1} and \cite[Proposition 2.6]{Han20} that for any space $X$, there exists a canonical isomorphism
\[
 \Omega^{\SO, \mathscr{Q}}_*(X) \otimes \Z\big[\tfrac{1}{2}\big] \cong H_*\big(X; \Z\big[\tfrac{1}{2}\big] \big) .
 \]
This implies that the image of $\beta$ in $\Omega^{\SO, \mathscr{Q}}_*( \sB \pi_1(V)) \otimes \Z\big[\tfrac{1}{2}\big]$ is equal to zero.

Thus there exists some $d \geq 0$ such that the image of $2^d \cdot \beta$ in $\Omega^{\SO, \mathscr{Q}}_*( \sB \pi_1(V))$ is equal to $0$.
Let $s, t \in \Z$ be integers with $s \cdot k + t \cdot 2^d = 1$.
This implies
\[
 s k \cdot \beta = \beta \in \Omega^{\SO, \mathscr{Q}}_*( \sB \pi_1(V)) .
\]
Let $W \to \sB \pi_1(V)$ be a compact $\mathscr{Q}$-bordism between $s k \cdot \beta $ and $\beta$.
More precisely, $W$ is a compact oriented $\mathscr{Q}_n$-manifold for some $n \geq 0$ (see \cite[p.\ 503]{Han20}) and $W \to \sB \pi_1(V)$ compatible with the singularity structure of~$W$ (see \cite[Definition 2.3]{Han20}).
Furthermore, along the boundary $\partial_0 V$ of~$W$ with its induced orientation (see \cite[Definition 2.1]{Han20}), this map restricts to the disjoint union of $s$ copies of $M \to \sB \pi_1(V)$, of $\varphi \colon V^{\rm op}\to \sB \pi_1(V)$ and of $f \colon V' \to \sB \pi_1(V)$.
Here, $V^{\rm op}$ is the manifold $V$ with the opposite orientation.

By applying surgeries in the interior of $W$, we can assume that the inclusion $V \hookrightarrow W$ is a~$2$-equivalence, i.e., it induces bijections on $\pi_0$ and $\pi_1$ and a~surjection on $\pi_2$.
This follows by a~similar argument as in the smooth case, see \cite[Theorem 1.5]{R-S1995}.\footnote{For some details one may consult the proof of \cite[Proposition 3.1]{SchickZenobi}.}

Now, since $M$ and $V'$ are \yes, and $\mathscr{Q}$ is a positive family of singularity types, the $\mathscr{Q}$-bordism principle \cite[Proposition 3.11]{Han20} shows that the manifold $V$ is \yes, completing the proof of Theorem \ref{prop:nonspin}.
\end{proof}

\begin{Remark} Theorem~\ref{prop:nonspin} together with a short proof has already appeared in \cite[Theorem~4.11]{RS}.
However, this discussion is incomplete if $H_*(\pi_1(V); \Z)$ contains $2$-torsion, which may produce Tor terms in the K\"unneth formula for the $E^2$-terms \smash{$H_p\big( \sB \pi_1(V) ; \Omega^{\SO}_q\big)_{(2)}$} in the Atiyah--Hirzebruch spectral sequence for \smash{$\Omega^{\SO}_*( \sB \pi_1(V))_{(2)}$}.
This is covered by Assertion~\ref{Assertion2} in the proof of our Proposition \ref{crucial}.
\end{Remark}

\subsection*{Acknowledgements}

This paper is the result of conversations in New York in the fall of 2018 and 2022.
B.H.~acknowledges support from NYU, the IAS Princeton and the DFG-funded Special Priority Program 2026 {\em Geometry at Infinity}.
We are grateful to the referees for their many useful comments, which helped to improve the manuscript.


\begin{thebibliography}{99}
\footnotesize\itemsep=0pt

\bibitem{Adem-Milgram}
Adem A., Milgram R.J., Cohomology of finite groups, 2nd ed., \textit{Grundlehren Math.
 Wiss.}, Vol.~309, \href{https://doi.org/10.1007/978-3-662-06280-7}{Springer}, Berlin, 2004.

\bibitem{baas1}
Baas N.A., On bordism theory of manifolds with singularities, \href{https://doi.org/10.7146/math.scand.a-11491}{\textit{Math.
 Scand.}} \textbf{33} (1973), 279--302.

\bibitem{BH10}
Brunnbauer M., Hanke B., Large and small group homology, \href{https://doi.org/10.1112/jtopol/jtq014}{\textit{J.~Topol.}}
 \textbf{3} (2010), 463--486, \href{https://arxiv.org/abs/0902.0869}{arXiv:0902.0869}.

\bibitem{C-E1956}
Cartan H., Eilenberg S., Homological algebra, Princeton University Press,
 Princeton, NJ, 1956.

\bibitem{Ce2020}
Cecchini S., A long neck principle for {R}iemannian spin manifolds with
 positive scalar curvature, \href{https://doi.org/10.1007/s00039-020-00545-1}{\textit{Geom. Funct. Anal.}} \textbf{30} (2020),
 1183--1223, \href{https://arxiv.org/abs/2002.07131}{arXiv:2002.07131}.

\bibitem{Ch-Li}
Chodosh O., Li C., Liokumovich Y., Classifying sufficiently connected {PSC}
 manifolds in~4 and~5 dimensions,  \href{https://doi.org/10.2140/gt.2023.27.1635}{\textit{Geom. Topol.}} \textbf{27} (2023),
 1635--1655, \href{https://arxiv.org/abs/2105.07306}{arXiv:2105.07306}.

\bibitem{CMS}
Chodosh O., Mantoulidis C., Schulze F., Improved generic regularity of
 codimension-1 minimizing integral currents, \textit{Ars Inven. Anal.}
 (2024), 3, 16~pages, \href{https://arxiv.org/abs/2306.13191}{arXiv:2306.13191}.

\bibitem{CF}
Conner P.E., Floyd E.E., Differentiable periodic maps, \textit{Ergeb. Math.
 Grenzgeb.}, Vol.~33, Springer, Berlin, 1964.

\bibitem{Fuehring}
F\"uhring S., A smooth variation of {B}aas--{S}ullivan theory and positive
 scalar curvature, \href{https://doi.org/10.1007/s00209-012-1105-x}{\textit{Math.~Z.}} \textbf{274} (2013), 1029--1046,
 \href{https://arxiv.org/abs/1203.4370}{arXiv:1203.4370}.

\bibitem{Fuehring_bordism}
F\"uhring S., Bordism and projective space bundles, \href{https://doi.org/10.17879/21089690084}{\textit{M\"unster~J. Math.}}
 \textbf{15} (2022), 355--388, \href{https://arxiv.org/abs/2006.15394}{arXiv:2006.15394}.

\bibitem{GdlH}
Ghys E., de~la Harpe P. (Editors), Sur les groupes hyperboliques d'apr\`es
 {M}ikhael {G}romov, \textit{Progr. Math.}, Vol.~83, \href{https://doi.org/10.1007/978-1-4684-9167-8}{Birkh\"auser}, Boston, MA,
 1990.

\bibitem{Gr1996}
Gromov M., Positive curvature, macroscopic dimension, spectral gaps and higher
 signatures, in Functional {A}nalysis on the {E}ve of the 21st {C}entury,
 {V}ol.~{II} ({N}ew {B}runswick, {NJ}, 1993), \textit{Progr. Math.}, Vol. 132,
  \href{https://doi.org/10.1007/978-1-4612-4098-3_1}{Birkh\"auser}, Boston, MA, 1996, 1--213.

\bibitem{Gr2018}
Gromov M., Metric inequalities with scalar curvature, \href{https://doi.org/10.1007/s00039-018-0453-z}{\textit{Geom. Funct.
 Anal.}} \textbf{28} (2018), 645--726, \href{https://arxiv.org/abs/1710.04655}{arXiv:1710.04655}.

\bibitem{Gr21}
Gromov M., Four lectures on scalar curvature, in Perspectives in {S}calar
 {C}urvature. {V}ol.~1, \href{https://doi.org/10.1142/97898112732230001}{World Scientific Publishing}, Hackensack, NJ, 2023,
 1--514, \href{https://arxiv.org/abs/1908.10612}{arXiv:1908.10612}.

\bibitem{GL}
Gromov M., Lawson Jr. H.B., The classification of simply connected manifolds of
 positive scalar curvature, \href{https://doi.org/10.2307/1971103}{\textit{Ann. of Math.}} \textbf{111} (1980),
 423--434.

\bibitem{G-L1983}
Gromov M., Lawson Jr. H.B., Positive scalar curvature and the {D}irac operator
 on complete {R}iemannian manifolds, \href{https://doi.org/10.1007/BF02953774}{\textit{Inst. Hautes \'Etudes Sci. Publ.
 Math.}} (1983), 83--196.

\bibitem{G-X-Y2020}
Guo H., Xie Z., Yu G., Quantitative {K}-theory, positive scalar curvature, and
 bandwidth, in Perspectives in scalar curvature. {V}ol.~2, \href{https://doi.org/10.1142/12644-vol2}{World Scientific
 Publishing}, Hackensack, NJ, 2023, 763--798, \href{https://arxiv.org/abs/2010.01749}{arXiv:2010.01749}.

\bibitem{Han20}
Hanke B., Positive scalar curvature on manifolds with odd order abelian
 fundamental groups, \href{https://doi.org/10.2140/gt.2021.25.497}{\textit{Geom. Topol.}} \textbf{25} (2021), 497--546,
 \href{https://arxiv.org/abs/1908.00944}{arXiv:1908.00944}.

\bibitem{Hitchin}
Hitchin N., Harmonic spinors, \href{https://doi.org/10.1016/0001-8708(74)90021-8}{\textit{Adv. Math.}} \textbf{14} (1974), 1--55.

\bibitem{Jo2004}
Joachim M., Toral classes and the {G}romov--{L}awson--{R}osenberg conjecture
 for elementary abelian 2-groups, \href{https://doi.org/10.1007/s00013-004-1080-5}{\textit{Arch. Math. (Basel)}} \textbf{83}
 (2004), 461--466.

\bibitem{K-M2012}
Kahn J., Markovic V., Immersing almost geodesic surfaces in a closed hyperbolic
 three manifold, \href{https://doi.org/10.4007/annals.2012.175.3.4}{\textit{Ann. of Math.}} \textbf{175} (2012), 1127--1190,
 \href{https://arxiv.org/abs/0910.5501}{arXiv:0910.5501}.

\bibitem{Loh2018}
Lohkamp J., Minimal smoothings of area minimizing cones, \href{https://arxiv.org/abs/1810.03157}{arXiv:1810.03157}.

\bibitem{Mar74}
Margolis H.R., Eilenberg--{M}ac {L}ane spectra, \href{https://doi.org/10.2307/2038905}{\textit{Proc. Amer. Math. Soc.}}
 \textbf{43} (1974), 409--415.

\bibitem{MOP08}
Matthey M., Oyono-Oyono H., Pitsch W., Homotopy invariance of higher signatures
 and 3-manifold groups, \href{https://doi.org/10.24033/bsmf.2547}{\textit{Bull. Soc. Math. France}} \textbf{136} (2008),
 1--25, \href{https://arxiv.org/abs/math.AT/0412363}{arXiv:math.AT/0412363}.

\bibitem{Milnor1962}
Milnor J., A unique decomposition theorem for 3-manifolds, \href{https://doi.org/10.2307/2372800}{\textit{Amer.~J.
 Math.}} \textbf{84} (1962), 1--7.

\bibitem{MP2021}
Morgan F., Pansu P., A list of open problems in differential geometry,
 \href{https://doi.org/10.1007/s40863-019-00141-8}{\textit{S\~ao Paulo~J. Math. Sci.}} \textbf{15} (2021), 305--321.

\bibitem{Raede22}
R\"ade D., Scalar and mean curvature comparison via {$\mu $}-bubbles,
 \href{https://doi.org/10.1007/s00526-023-02520-8}{\textit{Calc. Var. Partial Differential Equations}} \textbf{62} (2023), 187,
 39~pages, \href{https://arxiv.org/abs/2104.10120}{arXiv:2104.10120}.

\bibitem{Ros1986}
Rosenberg J., {$C^\ast$}-algebras, positive scalar curvature, and the {N}ovikov
 conjecture.~{III}, \href{https://doi.org/10.1016/0040-9383(86)90047-9}{\textit{Topology}} \textbf{25} (1986), 319--336.

\bibitem{R-S1995}
Rosenberg J., Stolz S., A ``stable'' version of the {G}romov--{L}awson
 conjecture, in The {\v{C}}ech {C}entennial, \textit{Contemp. Math.}, Vol.
 181, \href{https://doi.org/10.1090/conm/181/02046}{American Mathematical Society}, Providence, RI, 1995, 405--418,
 \href{https://arxiv.org/abs/dg-ga/9407002}{arXiv:dg-ga/9407002}.

\bibitem{RS}
Rosenberg J., Stolz S., Metrics of positive scalar curvature and connections
 with surgery, in Surveys on {S}urgery {T}heory, {V}ol.~2, \textit{Ann. of
 Math. Stud.}, Vol. 149, \href{https://doi.org/10.1515/9781400865215-010}{Princeton University Press}, Princeton, NJ, 2001,
 353--386.

\bibitem{Sch98}
Schick T., A counterexample to the (unstable) {G}romov--{L}awson--{R}osenberg
 conjecture, \href{https://doi.org/10.1016/S0040-9383(97)00082-7}{\textit{Topology}} \textbf{37} (1998), 1165--1168,
 \href{https://arxiv.org/abs/math/0403063}{arXiv:math/0403063}.

\bibitem{SchickZenobi}
Schick T., Zenobi V.F., Positive scalar curvature due to the cokernel of the
 classifying map, \href{https://doi.org/10.3842/SIGMA.2020.129}{\textit{SIGMA}} \textbf{16} (2020), 129, 12~pages,
 \href{https://arxiv.org/abs/2006.15965}{arXiv:2006.15965}.

\bibitem{S-Y1979}
Schoen R., Yau S.-T., Existence of incompressible minimal surfaces and the
 topology of three-dimensional manifolds with nonnegative scalar curvature,
 \href{https://doi.org/10.2307/1971247}{\textit{Ann. of Math.}} \textbf{110} (1979), 127--142.

\bibitem{SY}
Schoen R., Yau S.-T., On the structure of manifolds with positive scalar
 curvature, \href{https://doi.org/10.1007/BF01647970}{\textit{Manuscripta Math.}} \textbf{28} (1979), 159--183.

\bibitem{S-Y2017}
Schoen R., Yau S.-T., Positive scalar curvature and minimal hypersurface
 singularities, in Surveys in {D}ifferential {G}eometry, \textit{Surv. Differ.
 Geom.}, Vol.~24, \href{https://doi.org/10.4310/SDG.2019.v24.n1.a10}{Int. Press}, Boston, MA, 2022, 441--480, \href{https://arxiv.org/abs/1704.05490}{arXiv:1704.05490}.

\bibitem{Smale}
Smale N., Generic regularity of homologically area minimizing hypersurfaces in
 eight-dimensional manifolds, \href{https://doi.org/10.4310/CAG.1993.v1.n2.a2}{\textit{Comm. Anal. Geom.}} \textbf{1} (1993),
 217--228.

\bibitem{Stolz92}
Stolz S., Simply connected manifolds of positive scalar curvature, \href{https://doi.org/10.2307/2946598}{\textit{Ann.
 of Math.}} \textbf{136} (1992), 511--540.

\bibitem{Stolz_splitting}
Stolz S., Splitting certain {${\rm MSpin}$}-module spectra,
 \href{https://doi.org/10.1016/0040-9383(94)90040-X}{\textit{Topology}} \textbf{33} (1994), 159--180.

\bibitem{S-W-Z2021}
Su G., Wang X., Zhang W., Nonnegative scalar curvature and area decreasing maps
 on complete foliated manifolds, \href{https://doi.org/10.1515/crelle-2022-0038}{\textit{J.~Reine Angew. Math.}} \textbf{790}
 (2022), 85--113, \href{https://arxiv.org/abs/2104.03472}{arXiv:2104.03472}.

\bibitem{Wall60}
Wall C.T.C., Determination of the cobordism ring, \href{https://doi.org/10.2307/1970136}{\textit{Ann. of Math.}}
 \textbf{72} (1960), 292--311.

\bibitem{Ze2020}
Zeidler R., Width, largeness and index theory, \href{https://doi.org/10.3842/SIGMA.2020.127}{\textit{SIGMA}} \textbf{16}
 (2020), 127, 15~pages, \href{https://arxiv.org/abs/2008.13754}{arXiv:2008.13754}.

\bibitem{Ze2019}
Zeidler R., Band width estimates via the {D}irac operator,
 \href{https://doi.org/10.4310/jdg/1668186790}{\textit{J.~Differential Geom.}} \textbf{122} (2022), 155--183,
 \href{https://arxiv.org/abs/1905.08520}{arXiv:1905.08520}.

\end{thebibliography}
\end{document}
